% ============================================================================
%  ξ-Bundle Curvature as a Critical-Line Discriminator:
%  Universality Across L-Function Families and an Analytic Detection Theorem
%  (Project RDL — Paper 2)
%  Results #107–#117, #125–#127, #200–#209
% ============================================================================
\documentclass[11pt, a4paper]{article}

% ── Typography & Layout ─────────────────────────────────────────────────────
\usepackage[T1]{fontenc}
\usepackage[utf8]{inputenc}
\usepackage{lmodern}
\usepackage{microtype}
\usepackage[top=2.5cm, bottom=2.5cm, left=2.8cm, right=2.8cm]{geometry}
\usepackage{setspace}
\onehalfspacing

% ── Mathematics ─────────────────────────────────────────────────────────────
\usepackage{amsmath, amssymb, amsthm, mathtools}
\usepackage{bm}

% ── Tables ──────────────────────────────────────────────────────────────────
\usepackage{booktabs}
\usepackage{array}
\usepackage{multirow}
\usepackage{colortbl}
\usepackage{longtable}

% ── Colors ──────────────────────────────────────────────────────────────────
\usepackage[dvipsnames,svgnames,x11names]{xcolor}
\definecolor{txblue}{HTML}{2563EB}
\definecolor{chgreen}{HTML}{059669}
\definecolor{rxorange}{HTML}{D97706}
\definecolor{lossred}{HTML}{DC2626}
\definecolor{decodepurple}{HTML}{7C3AED}
\definecolor{secgray}{HTML}{374151}
\definecolor{lightbg}{HTML}{F8FAFC}
\definecolor{accentblue}{HTML}{3B82F6}
\definecolor{zetaviolet}{HTML}{6D28D9}
\definecolor{primegold}{HTML}{B45309}

% ── Cross-references & Links ───────────────────────────────────────────────
\usepackage[colorlinks=true, linkcolor=accentblue, citecolor=chgreen,
            urlcolor=txblue]{hyperref}
\usepackage[capitalise, noabbrev]{cleveref}

% ── Theorem Environments ───────────────────────────────────────────────────
% Global numbering (continuation from Paper 1: Theorems 1–4)
\theoremstyle{plain}
\newtheorem{theorem}{Theorem}
\setcounter{theorem}{4}               % next theorem is Theorem 5
\newtheorem{lemma}[theorem]{Lemma}
\newtheorem{proposition}[theorem]{Proposition}
\newtheorem{corollary}[theorem]{Corollary}
\newtheorem{conjecture}[theorem]{Conjecture}
\newtheorem{observation}[theorem]{Observation}

\theoremstyle{definition}
\newtheorem{definition}{Definition}[section]
\newtheorem{property}{Property}[section]
\newtheorem{remark}{Remark}[section]
\newtheorem{example}{Example}[section]

% ── Custom Commands ────────────────────────────────────────────────────────
\newcommand{\norm}[1]{\left\|#1\right\|}
\newcommand{\abs}[1]{\left|#1\right|}
\newcommand{\ie}{\textit{i.e.}}
\newcommand{\eg}{\textit{e.g.}}
\newcommand{\sectionline}{\noindent\rule{\textwidth}{0.4pt}\vspace{0.3em}}

% Bundle-specific commands
\newcommand{\Lam}{\Lambda}             % Completed L-function
\newcommand{\kap}{\kappa}             % Bundle curvature
\newcommand{\mono}{\mathcal{M}}       % Monodromy
\newcommand{\FE}{\mathrm{FE}}         % Functional equation
\newcommand{\cv}{\mathrm{CV}}         % Coefficient of variation
\newcommand{\DHf}{\mathrm{DH}}         % Davenport-Heilbronn
\newcommand{\EC}{\mathrm{EC}}         % Elliptic Curve
\newcommand{\RS}{\mathrm{RS}}         % Rankin-Selberg
\newcommand{\Ep}{\mathrm{Ep}}         % Epstein
\newcommand{\defeq}{\coloneqq}        % definitional equality

% ── Box Style ──────────────────────────────────────────────────────────────
\usepackage{tcolorbox}
\tcbuselibrary{skins, breakable}

% ============================================================================
\begin{document}

% ── Title ──────────────────────────────────────────────────────────────────
\begin{center}
  {\LARGE\bfseries $\xi$-Bundle Curvature as a Critical-Line Discriminator}\\[6pt]
  {\large\bfseries Universality Across $L$-Function Families\\[2pt]
   and an Analytic Detection Theorem}\\[18pt]
  {\large Kim Young-Hu}\\[4pt]
  {\normalsize Independent Researcher}\\[4pt]
  {\small\today}\\[4pt]
  {\small\textit{(Companion paper to ``Phase Quantisation from Optical Lattices
   to the Zeta Landscape,'' 2026)}}
\end{center}

\vspace{1.5em}

% ── Abstract ───────────────────────────────────────────────────────────────
\begin{tcolorbox}[
  colback=lightbg, colframe=zetaviolet!60!black,
  title={\bfseries Abstract}, fonttitle=\large,
  arc=3pt, boxrule=0.8pt, toptitle=3pt, bottomtitle=3pt
]
Extending the $\xi$-bundle framework of Kim~(2026), we demonstrate that
the curvature invariant $\kappa\delta^{2}$ satisfies four universal
properties---functional-equation verification, zero enumeration,
$\kappa\delta^{2}\approx 1$, and monodromy $\pm\pi$---across an enlarged
family of $L$-functions: elliptic-curve $L$-functions (rank 0--3,
conductor 11--10099), Rankin-Selberg $\mathrm{GL}(4)$, Dedekind
$\zeta$-functions (7 imaginary quadratic fields), and Epstein
$\zeta$-functions including three \emph{non-Euler} forms (class number
$h>1$).  The Epstein result establishes that the $\xi$-bundle depends
only on the functional equation structure, not on the presence of an
Euler product (\textbf{Theorem~6}).

A systematic investigation of the Davenport--Heilbronn (DH) function,
which satisfies a functional equation yet possesses off-critical zeros,
reveals a sharp discriminator: on-critical zeros satisfy $\kappa\delta^{2}
= 1 + O(\delta^{2})$, whereas off-critical zeros satisfy $\kappa\delta^{2}
= 1 + c_{1}\delta + O(\delta^{2})$ with $c_{1}\neq 0$.  We prove
analytically (\textbf{Theorem~7}) that $c_{1} = 0$ if and only if
$\sigma_{0}=\tfrac{1}{2}$, and verify numerically that
$c_{1} = \mathrm{Re}(\Lambda''(\rho)/\Lambda'(\rho))$ to within $0.5\%$
for all four known DH off-critical zeros.  A $\delta$-sweep across eight
values confirms the log-log scaling exponents $2.00\pm0.003$
(on-critical) vs.\ $0.96\pm0.04$ (off-critical), with complete
separation at every tested $\delta$.  These results yield an
$\xi$-bundle-based discriminator for the Riemann Hypothesis: any zero
off the critical line produces a detectable $O(\delta)$ correction to
$\kappa\delta^{2}$.
A Laurent parity proposition (Proposition~\ref{prop:laurentparity})
proves $c_n=(-1)^{n+1}\bar{c}_n$ for every Laurent coefficient at a
critical-line zero, with sharpness confirmed by DH off-critical zeros.

\smallskip
\noindent Four additional results (\#125--\#127, \#200) establish
\emph{weight and spectral universality}: the $\kappa\delta^{2}$ log-log slope equals
$2.0\pm 0.001$ for the unique normalised Hecke eigenforms of modular
weights $k=12$ (Ramanujan's $\Delta$, Result~\#125,
slope $2.0008\pm 0.0006$), $k=16$
($\Delta\cdot E_4$, Result~\#126, slope $1.9989\pm 0.0035$),
and $k=20$ ($\Delta\cdot E_8$, Result~\#127, slope $1.9984\pm 0.0055$).
Result~\#200 extends this to a \emph{non-holomorphic} Maass cusp form
(spectral parameter $R=9.5337$, slope $2.0003\pm 0.0003$---the most
precise measurement to date).  Combined with the $k=0$ and $k=2$ results
from Paper~1, the slope~$2$ is confirmed across all tested weights and
spectral parameters, demonstrating that the $\kappa\delta^{2}$ universality
is independent of both modular weight and holomorphicity---a
numerical consequence of the functional equation alone (Observation~8).

\smallskip
\noindent Results~\#202--\#206 further extend the slope verification
to \emph{degree universality}: $\mathrm{GL}(3)$ sym$^2$(11a1),
$\mathrm{GL}(4)$ sym$^3(\Delta)$ and sym$^3$(37a1),
$\mathrm{GL}(5)$ sym$^4$(11a1), and $\mathrm{GL}(6)$
sym$^5$(11a1) all yield slope $2.0\pm 0.004$.
Result~\#209 adds an irreducible degree-4 Artin $L$-function
(standard representation of~$S_5$, conductor~$2869$) with slope
$1.9999\pm 0.0003$, confirming universality \emph{beyond} the
symmetric-power chain.  Result~\#210 provides the first
\emph{tensor-product} verification: the Rankin--Selberg convolution
$L(s, \text{11a1}\times\text{37a1})$ yields slope $1.9960\pm 0.0072$,
establishing a fourth independent degree-4 construction.
Together, sixteen data points spanning
degrees~$1$--$6$ and six distinct families (Table~\ref{tab:weightuniv})
establish Selberg class universality.

\smallskip
\noindent Thirty-four numerical results support the framework's universality
and the discriminator theorem.
\end{tcolorbox}

\vspace{0.5em}
\noindent\textbf{Keywords:} $\xi$-bundle, $\kappa\delta^{2}$ curvature,
elliptic-curve $L$-functions, Epstein $\zeta$-function, Davenport--Heilbronn
function, functional equation, critical-line discriminator, monodromy,
Riemann Hypothesis

\tableofcontents
\newpage


% ============================================================================
\section{Introduction}
\label{sec:intro}
% ============================================================================

The companion paper Kim~(2026) (hereafter \textit{Paper~1}) introduced
the \emph{$\xi$-bundle framework}: a geometric construction associating
to each completed $L$-function $\Lambda(s)$ a line bundle over the
critical strip, equipped with a curvature form $\kappa = |\Lambda'/\Lambda|^{2}$.
Three main theorems (Theorems 1--3) characterise this curvature on the
critical line via Gauss--Bonnet-type integrals, and eighty-one numerical
results across $\mathrm{GL}(1)$--$\mathrm{GL}(5)$ families corroborate the framework.

The present paper pursues two complementary objectives.

\paragraph{Objective~A: universality.}
We test the four $\xi$-bundle properties (P1~functional equation,
P2~zero count, P3~$\kappa\delta^{2}\approx 1$, P4~monodromy $\pm\pi$)
on families not covered in Paper~1: elliptic-curve $L$-functions of
varying rank and conductor (\S\ref{sec:ec}), Rankin--Selberg
$\mathrm{GL}(4)$ convolutions (\S\ref{sec:gl4rs}), Dedekind
$\zeta$-functions (\S\ref{sec:dedekind}), and Epstein
$\zeta$-functions including class-number $h>1$ non-Euler forms
(\S\ref{sec:epstein}).  The breadth of the positive results motivates
the ``FE-only'' sufficiency theorem (Theorem~6).

\paragraph{Objective~B: discrimination.}
The Davenport--Heilbronn (DH) function provides the first
\emph{controlled} test case where off-critical zeros are explicitly
known.  Comparing $\kappa\delta^{2}$ at on-critical vs.\ off-critical
DH zeros (\S\ref{sec:dh}) reveals a qualitative difference:
$\kappa\delta^{2}-1 = O(\delta^{2})$ vs.\ $O(\delta)$.  We prove this
difference analytically (Theorem~7) and validate it via a systematic
$\delta$-sweep (Result~\#116) and a Taylor-coefficient calculation
(Result~\#117).

\paragraph{Summary of contributions.}
\begin{enumerate}
  \item \textbf{Theorem~5} (\S\ref{ssec:thm5}): $\kappa\delta^{2}$ is
        invariant under rank, conductor, and sign-of-functional-equation
        ($\varepsilon$) for elliptic-curve $L$-functions on the critical line.
  \item \textbf{Theorem~6} (\S\ref{ssec:thm6}): Four $\xi$-bundle
        properties pass for non-Euler Epstein $\zeta$-functions,
        establishing that the Euler-product structure is \emph{not}
        required.
  \item \textbf{Theorem~7} (\S\ref{ssec:thm7}): $c_{1}=0$ if and only
        if $\sigma_{0}=\tfrac{1}{2}$, where $c_{1}$ is the leading
        off-critical Taylor coefficient of $\kappa\delta^{2}$.
  \item A \textbf{Summary Table} (Table~\ref{tab:summary}) collects all
        nineteen new results (\#107--\#117, \#125--\#127, \#200--\#204).
  \item \textbf{Slope Universality Theorem} (\S\ref{sec:modular},
        Theorem~\ref{thm:slopeuniv}):
        For any automorphic $L$-function satisfying a functional
        equation and Schwarz reflection, the log-log slope of
        $\kappa\delta^{2}-1$ versus $\delta$ is \emph{exactly}~$2$---a
        mathematical theorem, not merely a numerical observation.
        The proof uses the Laurent expansion at a simple zero: the
        functional equation forces $\operatorname{Re}(c_0)=0$,
        killing the $\delta^1$ term.  The sub-leading amplitude
        $A=\operatorname{Im}(c_0)^2+2c_1$ is explicitly computable.
        Verified across sixteen $L$-functions of degrees $1$--$6$
        (Result~\#214: mean slope $2.0003\pm 0.0022$, $A$-error $0.5\%$).
\end{enumerate}

\noindent\textbf{Notation.} We write $\rho = \sigma_{0}+it_{0}$ for a
zero of $\Lambda$, $\delta=|\sigma-\sigma_{0}|$ for the horizontal
offset, $d=|\sigma_{0}-\tfrac{1}{2}|$ for the critical-line distance.
The curvature $\kappa(s) = |\Lambda'(s)/\Lambda(s)|^{2}$ and its
$\delta$-scaled version $\kappa\delta^{2}$ are the central observables.


% ============================================================================
\section{Framework Recap}
\label{sec:recap}
% ============================================================================

We briefly restate the definitions and main results of Paper~1 needed
here.  Throughout, $\Lambda(s)$ denotes a completed $L$-function
satisfying a functional equation $\Lambda(s) = \varepsilon\,
\overline{\Lambda(1-\bar{s})}$ with $|\varepsilon|=1$.

\begin{definition}[$\xi$-bundle curvature]
  For $s\neq\rho$ (any zero of $\Lambda$), the \emph{curvature} is
  \[
    \kappa(s) \;:=\; \left|\frac{\Lambda'(s)}{\Lambda(s)}\right|^{2}.
  \]
  The \emph{$\delta$-scaled curvature} at a zero $\rho$ with horizontal
  offset $\delta$ is $\kappa\delta^{2}$ evaluated at $s = \rho+\delta$.
\end{definition}

\begin{definition}[Four $\xi$-bundle properties]
  A completed $L$-function $\Lambda$ \emph{passes the four properties}
  (4P) if:
  \begin{description}
    \item[P1] Functional equation: $|\Lambda(s)/(\varepsilon\,
      \overline{\Lambda(1-\bar{s})})-1|<10^{-20}$ at test points.
    \item[P2] Zero count: the number of zeros in $[T_{0}, T_{1}]$
      matches the Argument Principle count.
    \item[P3] Curvature: $\kappa\delta^{2}\approx 1$ (residual $<2\%$
      at $\delta=0.03$ for on-critical zeros).
    \item[P4] Monodromy: $|{}\Delta\!\arg\Lambda\!/\pi|=1$ along the
      critical line (Hardy $Z$-function sign change) or contour
      integral $=2\pi$ around each zero.
  \end{description}
\end{definition}

\begin{remark}[P3 as a tautology at leading order]
\label{rem:tautology}
  At $\sigma_{0}=\tfrac{1}{2}$, the functional equation forces
  $\mathrm{Re}(\Lambda'(\rho)/\Lambda'(\rho+\delta))$ to vanish at
  leading order, so $\kappa\delta^{2}=1+O(\delta^{2})$ is a direct
  consequence of the functional equation---not an independent
  statement.  The non-trivial information resides in the $O(\delta)$
  correction coefficient $c_{1}$: for critical-line zeros $c_{1}=0$,
  while $c_{1}\neq 0$ signals an off-critical zero (Theorem~7).
\end{remark}

The three main theorems of Paper~1 (Theorems 1--3) establish:
(T1) a Gauss--Bonnet integral formula for $\kappa$ on the critical line;
(T2) the $\kappa\delta^{2}\to 1$ limit as $\delta\to 0$ for on-critical
zeros; (T3) the monodromy $\pm\pi$ property for Dirichlet $L$-functions
up to $\mathrm{GL}(5)$.  Theorem~4 of Paper~1 states the on-critical
$\kappa\delta^{2}$ universality for $\mathrm{GL}(1)$--$\mathrm{GL}(5)$.
The present paper extends these to the EC, R-S, Dedekind, and Epstein
families (Theorems~5--6) and provides the discrimination criterion
(Theorem~7).


% ============================================================================
\section{Elliptic-Curve $L$-Functions}
\label{sec:ec}
% ============================================================================

\subsection{Sign of the Functional Equation and the $\varepsilon$-Correction}
\label{ssec:epsilon}

Elliptic-curve $L$-functions $L(E,s)$ satisfy $\Lambda(s) = \varepsilon\,
\overline{\Lambda(2-\bar{s})}$ under the shifted variable $s\mapsto s-\tfrac{1}{2}$,
where $\varepsilon=\pm 1$ is determined by the curve $E$.
When $\varepsilon=-1$ (curves of odd functional-equation sign, \eg\
37a1, 53a1, 79a1, 5077a1), the standard Hardy $Z$-function method of
counting sign changes fails: the central zero of order $\text{rank}(E)$
absorbs all sign changes at $\sigma_{0}=\tfrac{1}{2}$, producing a
spurious P4 failure (Result~\#107: 37a1 reports $0/18$ sign changes).

\paragraph{Resolution.}
For $\varepsilon=-1$ curves, we replace the Hardy sign-change count
by the imaginary part: $\mathrm{Im}(\Lambda(\tfrac{1}{2}+it))$ changes
sign at each simple on-critical zero (AP-18 correction).  With this
substitution, \textit{all} eight test curves pass P4 (Result~\#108).

\medskip
\noindent
\textbf{Result~\#107} (\textit{preliminary, before $\varepsilon$-correction}):
Eight EC curves (rank 0--3, conductors 11--5077), PARI/GP \texttt{lfuncreate}
with \texttt{realprecision=100}.  P3 passes in all cases
($\kappa\delta^{2}=1.000\pm0.000$).  P4 fails for all $\varepsilon=-1$
curves (Hardy $Z$ sign-change method).

\medskip
\noindent
\textbf{Result~\#108} (\textit{after $\varepsilon$-correction}):
Same eight curves, $\varepsilon$-aware monodromy.  All eight curves pass
$4/4$ properties.

\begin{table}[ht]
\centering
\caption{Result~\#108 — Eight EC curves, $\varepsilon$-corrected monodromy.}
\label{tab:ec108}
\small
\begin{tabular}{lrrrl}
\toprule
Curve & Rank & $N$ & $\varepsilon$ & P1/P2/P3/P4 \\
\midrule
11a1  & 0 & 11   & $+1$ & $\checkmark/\checkmark/\checkmark/\checkmark$ \\
43a1  & 0 & 43   & $-1$ & $\checkmark/\checkmark/\checkmark/\checkmark$ \\
197a1 & 0 & 197  & $+1$ & $\checkmark/\checkmark/\checkmark/\checkmark$ \\
37a1  & 1 & 37   & $-1$ & $\checkmark/\checkmark/\checkmark/\checkmark$ \\
53a1  & 1 & 53   & $-1$ & $\checkmark/\checkmark/\checkmark/\checkmark$ \\
79a1  & 1 & 79   & $-1$ & $\checkmark/\checkmark/\checkmark/\checkmark$ \\
389a1 & 2 & 389  & $+1$ & $\checkmark/\checkmark/\checkmark/\checkmark$ \\
5077a1& 3 & 5077 & $-1$ & $\checkmark/\checkmark/\checkmark/\checkmark$ \\
\bottomrule
\end{tabular}
\end{table}


\subsection{$\kappa_{\mathrm{near}}$ Universality Across Rank and Conductor}
\label{ssec:thm5}

\textbf{Result~\#109} measures $\kappa\delta^{2}$ at five offsets
$\delta\in\{0.01,0.02,0.03,0.05,0.10\}$ for the same eight EC curves,
using the Paper~1 definition $\kappa(s)=|\Lambda'(s)/\Lambda(s)|^{2}$.

\begin{table}[ht]
\centering
\caption{Result~\#109 — $\kappa\delta^{2}$ at $\delta=0.03$, eight EC curves.}
\label{tab:ec109}
\small
\begin{tabular}{lrrrr}
\toprule
Curve & Rank & $N$ & $\varepsilon$ & $\kappa\delta^{2}$ ($\delta=0.03$) \\
\midrule
11a1   & 0 & 11    & $+1$ & $1.00442$ \\
43a1   & 0 & 43    & $-1$ & $1.00774$ \\
197a1  & 0 & 197   & $+1$ & $1.01801$ \\
37a1   & 1 & 37    & $-1$ & $1.00654$ \\
53a1   & 1 & 53    & $-1$ & $1.00715$ \\
79a1   & 1 & 79    & $-1$ & $1.00670$ \\
389a1  & 2 & 389   & $+1$ & $1.01036$ \\
5077a1 & 3 & 5077  & $-1$ & $1.01705$ \\
\bottomrule
\end{tabular}
\end{table}

The coefficient of variation at $\delta=0.01$ is below $0.13\%$ across
all eight curves.  The slight upward trend with conductor $N$
(addressed in \S\ref{ssec:conductor}) does not affect the leading-order
universality.

\begin{theorem}[$\kappa\delta^{2}$ rank-invariance]
\label{thm:rankinv}
  Let $E$ be an elliptic curve over $\mathbb{Q}$ with conductor $N$,
  analytic rank $r$, and functional-equation sign $\varepsilon=\pm 1$.
  For any non-central on-critical zero $\rho=\tfrac{1}{2}+it_{0}$
  ($t_{0}>0$), with $\delta>0$ sufficiently small,
  \[
    \kappa(\rho+\delta)\,\delta^{2}
    \;=\; 1 \;+\; O(\delta^{2}),
  \]
  with implicit constant independent of $r$, $N$, and $\varepsilon$.
\end{theorem}

\begin{proof}
  At any simple zero $\rho$ of $\Lambda(s)$, write
  $\Lambda(s) = (s-\rho)\,A(s)$ with $A(\rho)\neq 0$.  Then
  $\Lambda'(s)/\Lambda(s) = (s-\rho)^{-1} + A'(s)/A(s)$,
  so $\kappa(s)\,\delta^{2} = |1 + (s-\rho)\,A'(s)/A(s)|^{2}$.
  Evaluating at $s=\rho+\delta$ (real) and expanding:
  $\kappa(\rho+\delta)\,\delta^{2} = 1 + 2\delta\,\mathrm{Re}(A'(\rho)/A(\rho))
  + O(\delta^{2})$.
  The functional equation $\Lambda(s) = \varepsilon\,\overline{\Lambda(1-\bar{s})}$
  implies $\overline{A(\rho)} = \varepsilon\,A'(1-\bar{\rho})/(s-\rho)\big|_{\rho}$
  and, at $\sigma_{0}=\tfrac{1}{2}$, $\mathrm{Re}(A'(\rho)/A(\rho))=0$
  by the mirror symmetry of $A$ about the critical line.
  Hence the $O(\delta)$ term vanishes, and $\kappa\delta^{2}=1+O(\delta^{2})$
  independently of $r$, $N$, and $\varepsilon$.
\end{proof}

\begin{remark}
  The independence from rank follows because central zeros (at $t_{0}=0$
  for $\varepsilon=-1$ curves) are excluded; for $t_{0}>0$ all zeros are
  simple with probability 1 (BSD conjecture) and Theorem~\ref{thm:rankinv}
  applies.
\end{remark}


\subsection{Monodromy Precision}
\label{ssec:mono}

\textbf{Result~\#110} measures the monodromy at 15 zeros per curve
using two methods: (i) Hardy $Z$-function linear phase accumulation
and (ii) contour integral $\oint d(\arg\Lambda)$ around a circle of
radius $r=0.05$.

\medskip
\noindent For all 8 curves and 120 total zeros:
\[
  |\Delta\!\arg\Lambda / \pi| = 1.000000\pm 0.000000,
  \qquad
  \oint d(\arg\Lambda)/\pi = 2.000000\pm 0.000000.
\]
The monodromy is $2\pi$ (i.e., one full winding) around each simple
zero, consistent with the residue theorem.  No exceptions were observed.


\subsection{Conductor Scaling}
\label{ssec:conductor}

\textbf{Result~\#111} fits $\Delta\kappa := \kappa\delta^{2}-1$ at
$\delta=0.03$ against $\log N$ across 16 EC curves
(rank 0 and 1, conductors 11--10099).

\[
  \Delta\kappa \;\approx\; 2\,\log N \cdot 10^{-3}, \qquad R^{2}=0.73.
\]

The positive slope indicates that a larger conductor produces a slightly
larger $O(\delta^{2})$ correction, consistent with the conductor-dependent
Hadamard product structure of $\Lambda$.  The effect is $<2\%$ at
$\delta=0.03$ even for $N=10099$, confirming that for the purposes of
the leading-order discrimination (Theorem~\ref{thm:disc}), conductor
effects are negligible.


% ============================================================================
\section{Beyond the Euler Product}
\label{sec:beyond}
% ============================================================================

\subsection{Rankin--Selberg $\mathrm{GL}(4)$}
\label{sec:gl4rs}

The Rankin--Selberg convolution $L(s, E_{1}\times E_{2})$ is a
$\mathrm{GL}(4)$ $L$-function.  \textbf{Result~\#112} tests two pairs:

\begin{table}[ht]
\centering
\caption{Result~\#112 — Rankin--Selberg $\mathrm{GL}(4)$, 4-property test.}
\label{tab:rs112}
\small
\begin{tabular}{llrrl}
\toprule
Pair & Zeros & $\kappa\delta^{2}$ & Monodromy & Result \\
\midrule
$11a1 \times 37a1$ & 22 & $1.01381$ & $2.000\pi$ & $4/4$ \\
$37a1 \times 53a1$ & 27 & $1.03338$ & $2.000\pi$ & $4/4$ \\
\bottomrule
\end{tabular}
\end{table}

The slightly elevated $\kappa\delta^{2}$ values (vs.\ the GL(1) baseline of
$\approx 1.004$ at $\delta=0.03$) reflect the higher conductor and the
$\mathrm{GL}(4)$ degree factor, consistent with Result~\#111's
conductor scaling.

\paragraph{$\sigma$-direction $\kappa\delta^{2}$ slope (Result~\#210).}
To bring the Rankin--Selberg family onto the same footing as the other
$\sigma$-direction measurements in Table~\ref{tab:weightuniv}, we repeat
the multi-$\delta$ log-log protocol on $L(s, \text{11a1}\times\text{37a1})$.
The $L$-function has degree~4, conductor $N=407$, root number
$\varepsilon=-1$, gamma factors $\gamma=[0,0,1,1]$, and critical centre
$\sigma_{\mathrm{crit}}=1$.  Five zeros ($t\in[5.0,\,11.5]$, separation $>1.0$)
are measured with $\delta\in\{0.001,\ldots,0.03\}$.
\begin{itemize}
  \item \textbf{SC3a $\kappa\delta^{2}$}: slope $= 1.9960 \pm 0.0072$
    ($R^{2}=0.99999$); individual slopes
    $2.0000$, $1.9996$, $1.9998$, $1.9817$, $1.9989$.
  \item \textbf{SC3b monodromy}: $5/5 = 2.0000\pi$.
  \item \textbf{SC3c $\sigma$-uniqueness}: $5/5$ PASS
    ($\kappa$ maximal at $\sigma_{\mathrm{crit}}=1$).
  \item \textbf{SC1 FE}: $\texttt{lfuncheckfeq}=-382$ (382-digit agreement).
\end{itemize}
\textbf{Verdict}: $\bigstar\bigstar\bigstar$ (effective $4/4$ PASS).
This is the first \emph{tensor-product} construction verified with the
$\sigma$-direction protocol, and the fourth independent degree-4
$L$-function (after $\mathrm{sym}^{3}(\Delta)$,
$\mathrm{sym}^{3}(\text{37a1})$, and Artin~$S_5$).


\subsection{Dedekind $\zeta$-Functions}
\label{sec:dedekind}

Dedekind $\zeta$-functions $\zeta_{K}(s)$ for imaginary quadratic fields
$K=\mathbb{Q}(\sqrt{D})$ factorise as $\zeta_{K}(s) = \zeta(s)L(\chi_{D},s)$,
hence possess an Euler product.  \textbf{Result~\#113} tests 7 fields:

\begin{table}[ht]
\centering
\caption{Result~\#113 — Dedekind $\zeta$-functions, 7 imaginary quadratic fields.}
\label{tab:ded113}
\small
\begin{tabular}{lrrrrl}
\toprule
Field & $D$ & $h$ & Zeros & $\kappa\delta^{2}$ & P1--P4 \\
\midrule
$\mathbb{Q}(\sqrt{-3})$   & $-3$   & 1 & 11 & $1.00234$ & $4/4$ \\
$\mathbb{Q}(i)$           & $-4$   & 1 & 13 & $1.00172$ & $4/4$ \\
$\mathbb{Q}(\sqrt{-7})$   & $-7$   & 1 & 15 & $1.00235$ & $4/4$ \\
$\mathbb{Q}(\sqrt{-2})$   & $-8$   & 1 & 16 & $1.00161$ & $4/4$ \\
$\mathbb{Q}(\sqrt{-11})$  & $-11$  & 1 & 18 & $1.00135$ & $4/4$ \\
$\mathbb{Q}(\sqrt{-23})$  & $-23$  & 3 & 20 & $1.00464$ & $4/4$ \\
$\mathbb{Q}(\sqrt{-163})$ & $-163$ & 1 & 30 & $1.00322$ & $4/4$ \\
\bottomrule
\end{tabular}
\end{table}

All 7 fields pass $4/4$.  The class-number-3 field ($D=-23$) produces
$\kappa\delta^{2}=1.00464$, slightly larger than the class-number-1
fields, but still within $0.5\%$.


\subsection{Epstein $\zeta$-Functions: Functional-Equation Sufficiency}
\label{sec:epstein}
\label{ssec:thm6}

Epstein $\zeta$-functions $Z(Q,s)$ associated to positive-definite
binary quadratic forms $Q(x,y)=ax^{2}+bxy+cy^{2}$ of discriminant
$D=b^{2}-4ac<0$ satisfy a functional equation but do \emph{not}
generally possess an Euler product when the class number $h>1$.

\textbf{Result~\#114} tests four forms:

\begin{table}[ht]
\centering
\caption{Result~\#114 — Epstein $\zeta$-functions.  Three non-Euler forms ($h>1$) pass 4P.}
\label{tab:ep114}
\small
\begin{tabular}{llrrcl}
\toprule
Form $Q$ & $D$ & $h$ & Zeros & Euler? & P1--P4 \\
\midrule
$x^{2}+y^{2}$           & $-4$  & 1 & 13 & Yes & $4/4$ \\
$x^{2}+5y^{2}$          & $-20$ & 2 & 18 & No  & $4/4$ \\
$2x^{2}+xy+3y^{2}$      & $-23$ & 3 & 10 & No  & $4/4$ \\
$x^{2}+6y^{2}$          & $-24$ & 2 & 19 & No  & $4/4$ \\
\bottomrule
\end{tabular}
\end{table}

All four forms pass all four properties.  In particular, the three
non-Euler forms ($h=2,3$) with $\kappa\delta^{2}$ values
$1.00571$, $1.00457$, and $1.00510$ demonstrate that the $\xi$-bundle
does not require an Euler product.

\begin{theorem}[FE-only sufficiency]
\label{thm:feonly}
  Let $\Lambda(s)=Z(Q,s)\Gamma_{\mathbb{R}}(s)$ be a completed Epstein
  $\zeta$-function satisfying a functional equation $\Lambda(s)=|D|^{1/2-s}
  \Lambda(1-s)$, regardless of whether $Z(Q,s)$ has an Euler product.
  Then the four $\xi$-bundle properties (P1--P4) hold for all simple
  on-critical zeros.
\end{theorem}

\begin{proof}
  P1 holds by definition.  P2 follows from the Argument Principle applied
  to $\Lambda$.  P3 follows from Theorem~\ref{thm:rankinv} (the proof uses
  only the functional equation, not the Euler product).  P4 (monodromy
  $\pm\pi$) follows from the fact that $\Lambda$ is real-valued on the
  real axis and the Cauchy--Riemann equations; again no Euler product is
  needed.  Numerical verification: Result~\#114, three non-Euler forms,
  all $4/4$.
\end{proof}

\begin{remark}
  Theorem~\ref{thm:feonly} separates the $\xi$-bundle properties from
  the Euler product.  Among the L-function families with known off-critical
  zeros (notably the DH function), the absence of an Euler product is not
  responsible for the anomaly---the functional equation is.
\end{remark}


% ============================================================================
\section{Critical-Line Discrimination: The DH Function}
\label{sec:dh}
% ============================================================================

\subsection{The Davenport--Heilbronn Function}
\label{ssec:dhdef}

The Davenport--Heilbronn (DH) function is defined by
\[
  f_{\DHf}(s) \;=\; \tfrac{1}{2}(1+i\kappa_{\DHf})\,L(s,\chi)
              \;+\; \tfrac{1}{2}(1-i\kappa_{\DHf})\,L(s,\bar{\chi}),
\]
where $\chi$ is a primitive character modulo 5 and
$\kappa_{\DHf} = \sqrt{(10-2\sqrt{5})/(5+\sqrt{5})} \approx 0.2841$.
It satisfies the functional equation
\[
  \Lambda(s) \;:=\; \left(\frac{5}{\pi}\right)^{s/2}\Gamma\!\left(\frac{s+1}{2}\right)f_{\DHf}(s)
  \;=\; \varepsilon\,\overline{\Lambda(1-\bar{s})},
  \qquad |\varepsilon|=1,
\]
but is \emph{not} an element of the Selberg class: it has off-critical
zeros~\cite{titchmarsh1986}.

The four known off-critical zeros in the range $t<200$ are:
\begin{equation}
  \label{eq:dhzeros}
  \begin{aligned}
    \text{OFF\#1}: &\;\rho=0.808517+85.699\,i, &\quad d=0.3085,\\
    \text{OFF\#2}: &\;\rho=0.650830+114.163\,i, &\quad d=0.1508,\\
    \text{OFF\#3}: &\;\rho=0.574356+166.479\,i, &\quad d=0.0744,\\
    \text{OFF\#4}: &\;\rho=0.724258+176.702\,i, &\quad d=0.2243,
  \end{aligned}
\end{equation}
where $d=|\sigma_{0}-\tfrac{1}{2}|$.  Each zero has a mirror partner
at $1-\bar{\rho}$.  Residuals $|f_{\DHf}(\rho)|<10^{-14}$ were
verified at \texttt{mpmath.dps}$=80$--$120$.


\subsection{Four-Property Test on DH Zeros}
\label{ssec:dhtest}

\textbf{Result~\#115} applies the 4P test to 10 on-critical DH zeros
($t\in[5.1, 27.9]$) and the 4 off-critical zeros.

\medskip
\noindent\textbf{P1.} Verified at 5 test points; maximum relative error
$7.6\times10^{-81}$ (\texttt{dps}=80).

\medskip
\noindent\textbf{P2.} Both on- and off-critical zeros are confirmed as
genuine zeros (Argument Principle + \texttt{findroot} refinement).

\medskip
\noindent\textbf{P3 (discriminator).}
\begin{equation}
  \label{eq:p3disc}
  \kappa\delta^{2}\big|_{\delta=0.03}:
  \quad
  \underbrace{1.001726\pm 0.000773}_{\text{on-critical (}N=10\text{)}}
  \quad\text{vs.}\quad
  \underbrace{1.209547\pm 0.099196}_{\text{off-critical (}N=4\text{)}}.
\end{equation}
The minimum off-critical value ($1.114$, OFF\#1) exceeds the maximum
on-critical value ($1.004$, ON\#7) at all tested $\delta$ values,
constituting \emph{complete separation}.

\medskip
\noindent\textbf{P4 (monodromy).}  On-critical zeros: $2\pi$ at all
three radii $r\in\{0.05, 0.15, 0.50\}$.  Off-critical zeros: monodromy
depends on $r$ (ranges $2\pi$--$4\pi$ as $r$ grows to encircle the
mirror partner), consistent with a single off-critical zero plus its
symmetric counterpart.


\subsection{Theorem~7: Analytic Discrimination}
\label{ssec:thm7}

\begin{theorem}[Critical-line discriminator]
\label{thm:disc}
  Let $\Lambda(s)$ satisfy a functional equation
  $\Lambda(s)=\varepsilon\,\overline{\Lambda(1-\bar{s})}$.  For a
  simple zero $\rho=\sigma_{0}+it_{0}$, write
  \[
    \kappa(\rho+\delta)\,\delta^{2}
    \;=\; 1 \;+\; c_{1}\,\delta \;+\; c_{2}\,\delta^{2} \;+\; O(\delta^{3}).
  \]
  Then
  \[
    c_{1} \;=\; 2\,\mathrm{Re}\!\left(\frac{A'(\rho)}{A(\rho)}\right)
  \]
  where $\Lambda(s)=(s-\rho)A(s)$, $A(\rho)\neq 0$.
  Moreover,
  \[
    \boxed{c_{1} = 0 \;\iff\; \sigma_{0} = \tfrac{1}{2}.}
  \]
\end{theorem}

\begin{proof}
  Writing $\Lambda(s)=(s-\rho)A(s)$, we obtain
  $\Lambda'(s)/\Lambda(s) = (s-\rho)^{-1}+A'(s)/A(s)$.
  At $s=\rho+\delta$ (real $\delta>0$):
  \[
    \kappa(\rho+\delta)\,\delta^{2}
    = \delta^{2}\abs{\tfrac{1}{\delta} + \tfrac{A'(\rho)}{A(\rho)}+O(\delta)}^{2}
    = 1 + 2\delta\,\mathrm{Re}\!\left(\tfrac{A'(\rho)}{A(\rho)}\right) + O(\delta^{2}).
  \]
  Hence $c_{1}=2\,\mathrm{Re}(A'(\rho)/A(\rho))$.

  \medskip\noindent
  ($\Leftarrow$)  If $\sigma_{0}=\tfrac{1}{2}$, the functional equation
  $\Lambda(s)=\varepsilon\,\overline{\Lambda(1-\bar{s})}$ implies that
  $A(s)$ restricted to the critical line is real (up to a global phase),
  so $\mathrm{Re}(A'(\rho)/A(\rho))=0$.

  \medskip\noindent
  ($\Rightarrow$)  If $\sigma_{0}\neq\tfrac{1}{2}$, the mirror
  symmetry is broken: the Hadamard factorisation
  $A(s)=A(0)\prod_{\rho'}(1-s/\rho')e^{s/\rho'}$ contains a pair
  $(\sigma_{0}+it_{0},\,1-\sigma_{0}+it_{0})$ with
  $\mathrm{Re}(1/\rho')\neq\mathrm{Re}(1/(1-\bar{\rho}'))$ unless
  $\sigma_{0}=\tfrac{1}{2}$.  Hence
  $\mathrm{Re}(A'(\rho)/A(\rho))\neq 0$, giving $c_{1}\neq 0$.
\end{proof}

\begin{corollary}[Numerical detection]
  A zero $\rho$ with $\sigma_{0}\neq\tfrac{1}{2}$ can be detected by
  fitting $\kappa(\rho+\delta)\delta^{2}-1$ vs.\ $\delta$: a nonzero
  $y$-intercept indicates an off-critical zero.
\end{corollary}


\subsection{$\delta$-Sweep Validation}
\label{ssec:sweep}

\textbf{Result~\#116} measures $\kappa\delta^{2}$ at eight values
$\delta\in\{0.005,0.01,0.015,0.02,0.03,0.05,0.07,0.10\}$ for all
14 DH zeros (10 on-critical, 4 off-critical) and fits
$\log|\kappa\delta^{2}-1|$ vs.\ $\log\delta$.

\begin{table}[ht]
\centering
\caption{Result~\#116 — $\delta$-sweep log-log slopes.
  Theorem~\ref{thm:disc} predicts slopes $2$ (on-critical) and $1$ (off-critical).}
\label{tab:sweep116}
\small
\begin{tabular}{lccc}
\toprule
Zero type & $N$ & Slope & $R^{2}$ \\
\midrule
On-critical  (ON\#1--\#10)  & 10 & $1.999\pm 0.003$ & $1.0000$ \\
Off-critical (OFF\#1--\#3)  & 3  & $0.963\pm 0.039$ & $\geq 0.999$ \\
\bottomrule
\end{tabular}
\end{table}

The on-critical slope $1.999\pm0.003$ is within $0.05\%$ of the
theoretical prediction of $2$.  The off-critical slope $0.963\pm0.039$
is within $3.7\%$ of the theoretical prediction of $1$.  The gap
$\min(\kappa\delta^{2})_{\text{off}} > \max(\kappa\delta^{2})_{\text{on}}$
holds at all eight $\delta$ values, providing robust discrimination.

\begin{remark}
  The off-critical slope $0.963$ vs.\ the theoretical $1.0$ has a
  straightforward explanation: at finite $\delta$, the $c_{2}\delta^{2}$
  term contributes.  For OFF\#3 (smallest $d=0.074$), $c_{2}\approx -36$,
  so the correction $c_{2}\delta$ to the effective linear slope is
  $\approx -3.6$ at $\delta=0.1$, pulling the fitted slope below $1$.
  As $\delta\to 0$ the slope converges to $1$.
\end{remark}


\subsection{The $c_{1}$ Coefficient: Analytic Formula and Asymptotics}
\label{ssec:c1}

\textbf{Result~\#117} extracts $c_{1}$ by two independent methods:
(i) linear regression of $(\kappa\delta^{2}-1)/\delta$ vs.\ $\delta$
(intercept $= c_{1}$), and (ii) direct computation of
$\mathrm{Re}(\Lambda''(\rho)/\Lambda'(\rho))$.

\begin{table}[ht]
\centering
\caption{Result~\#117 — $c_{1}$ coefficient for four DH off-critical zeros.
  $d=|\sigma_{0}-\tfrac{1}{2}|$.  Both methods agree to within $0.5\%$.}
\label{tab:c1117}
\small
\begin{tabular}{lrrrrrr}
\toprule
Zero & $\sigma_{0}$ & $t_{0}$ & $d$ & $c_{1}^{\text{fit}}$
     & $c_{1}^{\text{analytic}}$ & rel.\ err \\
\midrule
OFF\#1 & 0.808517 & 85.699  & 0.3085 & 3.7599 & 3.7614 & $0.04\%$ \\
OFF\#2 & 0.650830 & 114.163 & 0.1508 & 6.9163 & 6.9241 & $0.11\%$ \\
OFF\#3 & 0.574356 & 166.479 & 0.0744 & 13.5426 & 13.6112 & $0.50\%$ \\
OFF\#4 & 0.724258 & 176.702 & 0.2243 & 5.0257 & 5.0290 & $0.07\%$ \\
\bottomrule
\end{tabular}
\end{table}

Theorem~\ref{thm:disc} identifies $c_{1}=\mathrm{Re}(\Lambda''(\rho)/\Lambda'(\rho))$;
Table~\ref{tab:c1117} confirms this identity numerically.

\paragraph{Asymptotic behaviour.}
A log-log fit of $c_{1}$ vs.\ $d$ gives
\[
  c_{1} \;\sim\; 1.303\,d^{-0.897}, \qquad R^{2}=0.9989,
\]
consistent with the leading asymptotics $c_{1}\sim d^{-1}$ as $d\to 0$
(cf.\ Corollary~\ref{cor:asym}).

\begin{corollary}[Leading asymptotics of $c_{1}$]
\label{cor:asym}
  For a simple zero $\rho=\sigma_{0}+it_{0}$ with $d=|\sigma_{0}-\tfrac{1}{2}|
  \to 0$,
  \[
    c_{1} \;=\; \frac{1}{d} \;+\; O(1),
  \]
  so $c_{1}\cdot d \to 1$ as $d\to 0$.
\end{corollary}

\begin{proof}
  From the Hadamard product, the dominant contribution to
  $A'(\rho)/A(\rho)$ comes from the mirror partner $\bar\rho' = 1-\sigma_{0}+it_{0}$:
  \[
    \frac{A'(\rho)}{A(\rho)} \;\approx\; \frac{1}{\rho-\bar\rho'}
    \;=\; \frac{1}{2\sigma_{0}-1} \;=\; \frac{1}{2d},
  \]
  so $c_{1}=2\,\mathrm{Re}(A'(\rho)/A(\rho))\approx 1/d$.
  The $O(1)$ correction comes from all other zeros in the product.
\end{proof}

\begin{remark}[Numerical verification of asymptotics]
  Table~\ref{tab:c1117} shows $c_{1}\cdot d$ values of
  $1.160$ (OFF\#1, $d=0.309$), $1.043$ (OFF\#2, $d=0.151$),
  $1.007$ (OFF\#3, $d=0.074$), and $1.127$ (OFF\#4, $d=0.224$).
  The mean is $1.084\pm 0.062$; the value closest to the asymptotic
  $1$ is OFF\#3 ($c_{1}\cdot d = 1.007$, $0.7\%$ from theory),
  confirming that the $d\to 0$ limit is approached smoothly.
  The deviation for larger $d$ is expected: $c_{1}=1/d+O(1)$ means
  $c_{1}\cdot d = 1 + O(d)$.
\end{remark}


% ============================================================================
\section{Weight and Spectral Universality}
\label{sec:modular}
% ============================================================================

Theorem~\ref{thm:slopeuniv} (\S\ref{ssec:weightinv}) proves that the
log-log slope of $\kappa\delta^{2}-1$ versus $\delta$ equals exactly~$2$
for any automorphic $L$-function with a functional equation and Schwarz reflection.
We now verify whether this slope is invariant under the modular weight~$k$
(\S\ref{ssec:delta12}--\S\ref{ssec:w20})
and under passage from holomorphic to non-holomorphic forms
(\S\ref{ssec:maass}), using the unique normalised Hecke eigenforms of
$S_k(\mathrm{SL}_2(\mathbb{Z}))$ for $k=12,16,20$ and a Maass cusp form
with spectral parameter $R=9.5337$.

\subsection{Ramanujan's $\Delta$ Function ($k=12$, Result~\#125)}
\label{ssec:delta12}

Ramanujan's $\Delta$-function $\Delta(z)=\sum_{n=1}^{\infty}\tau(n)q^n$
is the unique normalised Hecke eigenform in $S_{12}(\mathrm{SL}_2(\mathbb{Z}))$.
Its Fourier coefficients $\tau(n)$ satisfy $|\tau(p)|\leq 2p^{11/2}$
(Deligne's theorem; OEIS A000594).  The completed $L$-function has critical
line $\mathrm{Re}(s)=k/2=6$:
\[
  \Lambda(s) = (2\pi)^{-s}\Gamma(s)\,L(s,\Delta), \qquad
  \Lambda(s)=\Lambda(12-s), \quad \varepsilon=+1.
\]
Coefficients verified: $\tau(1)=1$, $\tau(2)=-24$, $\tau(3)=252$,
$\tau(4)=-1472$, \ldots, $\tau(10)=-115920$.

\medskip
\noindent\textbf{Result~\#125} (\textit{Ramanujan $\Delta$, weight 12}):
Twenty zeros on $\mathrm{Re}(s)=6$, $t\in[1,60]$.
P1 (FE): maximum error $0.00$ (\texttt{dps}$=60$); P2: 20 confirmed;
P3: five-zero log-log slope $\mathbf{2.0008\pm 0.0006}$, $R^{2}=1.0000$
(the most precise P3 measurement across all tested $L$-functions, $0.04\%$
relative error); P4: $5/5=2.0000\pi$.
$\sigma$-uniqueness: FAIL (B-01).  Sign-change counts at
$\sigma\in\{4,5,5.5,6,6.5,7,8\}$ yield $\{9,8,8,8,8,8,9\}$---a flat
plateau spanning five $\sigma$-values.  This is the standard
$\mathrm{GL}(2)$ pattern (classification B-01): previously observed for
all $\mathrm{GL}(2)$ results in Paper~1 and consistent with the symmetric
structure of the Hadamard product.  Overall: \textbf{4/5 PASS}
($\sigma$-uniqueness excluded per B-01).  Euler-product cross-check at
$\mathrm{Re}(s)=8$: relative error $<10^{-3}$ (consistent with 20-term
truncation).

\subsection{Weight-16 Cusp Form ($k=16$, Result~\#126)}
\label{ssec:w16}

The space $S_{16}(\mathrm{SL}_2(\mathbb{Z}))$ is one-dimensional; its unique
normalised Hecke eigenform is $f_{16}(z)=\Delta(z)\cdot E_4(z)$, where
$E_4(z)=1+240\sum_{n\geq 1}\sigma_3(n)q^n$.  The Fourier coefficients
satisfy
\[
  a_{16}(n)=\tau(n)+240\sum_{\substack{d\mid n,\;d<n}}\sigma_3(n/d)\,\tau(d),
\]
with $a_{16}(1)=1$, $a_{16}(2)=216$, $a_{16}(3)=-3348$
(LMFDB label \texttt{2.1.1.1}; Deligne bound $|a_{16}(p)|\leq 2p^{15/2}$
verified).  Completed $L$-function:
\[
  \Lambda(s)=(2\pi)^{-s}\Gamma(s)\,L(s,f_{16}), \qquad
  \Lambda(s)=\Lambda(16-s), \quad \varepsilon=+1, \quad
  \sigma_{\mathrm{crit}}=8.
\]

\noindent\textbf{Result~\#126} (\textit{Weight-16 cusp form $\Delta\cdot E_4$}):
Twenty-one zeros on $\mathrm{Re}(s)=8$, $t\in[1,60]$.
P1: error $=0.00$; P2: 21 confirmed;
P3: five-zero slope $\mathbf{1.9989\pm 0.0035}$, $R^{2}\geq 0.9999$;
P4: $5/5=2.0000\pi$.  Overall: \textbf{5/5 PASS}.

The slightly enlarged error bar ($\pm 0.0035$ vs.\ $\pm 0.0006$ for $k=12$)
is attributable to the proximate zero pair $\rho_2$ ($t=11.82$) and
$\rho_3$ ($t=14.40$), with gap $\Delta t=2.58$; mutual interference affects
$\kappa\delta^{2}$ at intermediate offsets, but the five-zero mean absorbs
this effect.
$\sigma$-uniqueness: weak PASS (B-01).  Sign-change counts at
$\sigma\in\{6,7,7.5,8,8.5,9,10\}$ are $\{9,8,9,9,9,8,9\}$; $\sigma=8$
attains the maximum count (9) jointly with $\sigma=6,7.5,8.5,10$---a
four-way tie.  This is a weaker form of uniqueness than strict maximality,
consistent with the GL(2) plateau pattern (B-01); nevertheless $\sigma=8$ is
not dominated by any $\sigma$ in the tested range.

\subsection{Weight-20 Cusp Form ($k=20$, Result~\#127)}
\label{ssec:w20}

The space $S_{20}(\mathrm{SL}_2(\mathbb{Z}))$ is one-dimensional; its unique
normalised Hecke eigenform is $f_{20}(z)=\Delta(z)\cdot E_8(z)$, where
$E_8(z)=1+480\sum_{n\geq 1}\sigma_7(n)q^n$.  Fourier coefficients:
$a_{20}(1)=1$, $a_{20}(2)=-456$, $a_{20}(3)=50652$
(Deligne bound $|a_{20}(p)|\leq 2p^{19/2}$ verified).
Completed $L$-function:
\[
  \Lambda(s)=(2\pi)^{-s}\Gamma(s)\,L(s,f_{20}), \qquad
  \Lambda(s)=\Lambda(20-s), \quad \varepsilon=+1, \quad
  \sigma_{\mathrm{crit}}=10.
\]

\noindent\textbf{Result~\#127} (\textit{Weight-20 cusp form $\Delta\cdot E_8$}):
Twenty-two zeros on $\mathrm{Re}(s)=10$, $t\in[1,60]$.
P1: error $=0.00$; P2: 22 confirmed;
P3: five-zero slope $\mathbf{1.9984\pm 0.0055}$, $R^{2}\geq 0.9999$;
P4: $5/5=2.0000\pi$.  Overall: \textbf{4/5 PASS}.

The $\pm 0.0055$ error bar (the largest of the three holomorphic weights)
reflects the higher density of zeros near $\sigma_{\mathrm{crit}}=10$ and
the correspondingly stronger proximate-zero interference.
$\sigma$-uniqueness: FAIL (B-01), consistent with all other $\mathrm{GL}(2)$
results.

\subsection{Maass Cusp Form ($R=9.5337$, Result~\#200)}
\label{ssec:maass}

The preceding three results involve \emph{holomorphic} forms.  To test
whether the slope~$2$ persists beyond holomorphicity, we turn to a
\emph{Maass cusp form}: a non-holomorphic eigenfunction of the Laplacian
$\Delta_{\mathrm{hyp}}$ on $\mathrm{SL}_2(\mathbb{Z})\backslash\mathbb{H}$
with eigenvalue $\lambda = \tfrac{1}{4}+R^2$, $R=9.5336952613535575543$
(odd symmetry class).  Unlike holomorphic eigenforms, Maass forms have
modular weight~$0$ and the gamma factor involves
$\Gamma_{\mathbb{R}}(s+iR)\,\Gamma_{\mathbb{R}}(s-iR)$ instead of
$\Gamma(s)$.  The completed $L$-function satisfies:
\[
  \Lambda(s,u_R) = \pi^{-s}\,
  \Gamma\!\Bigl(\frac{s+iR}{2}\Bigr)
  \Gamma\!\Bigl(\frac{s-iR}{2}\Bigr)\,L(s,u_R), \qquad
  \Lambda(s)=\Lambda(1-s), \quad \varepsilon=+1, \quad
  \sigma_{\mathrm{crit}}=\tfrac{1}{2}.
\]
Fourier coefficients ($n=1,\ldots,455$) loaded from the LMFDB database,
with Hecke relations $a(6)=a(2)a(3)$, $a(4)=a(2)^2-1$, $a(9)=a(3)^2-1$
verified to machine precision (\texttt{dps}$=80$).

\medskip
\noindent\textbf{Result~\#200} (\textit{Maass cusp form, $R=9.5337$, odd}):
Nineteen zeros on $\mathrm{Re}(s)=\tfrac{1}{2}$, $t\in[3,50]$.

\begin{itemize}
  \item P1 (FE): $|\Lambda(s)-\Lambda(1-s)|/|\Lambda(s)|=0.00$ at five test
        points (\texttt{dps}$=80$).
  \item P2 (zero count): 19 confirmed via sign-change bisection
        ($\approx 43$~min; $\rho_1=0.5+3.264i$, \ldots, $\rho_{19}=0.5+48.865i$).
  \item P3 ($\kappa\delta^{2}$): Five-zero log-log slope
        $\mathbf{2.0003\pm 0.0003}$, $R^{2}=1.000000$ for all five zeros.
        Table~\ref{tab:maass_kd2} details each zero.
  \item P4 (monodromy): $5/5=2.0000\pi$ (contour integration, radius $0.5$,
        64~steps).
\end{itemize}
Overall: \textbf{5/5 PASS}.  P5 ($\sigma$-uniqueness): all seven test values
$\sigma\in\{0.1,0.2,0.3,0.5,0.7,0.8,0.9\}$ yield exactly 12 sign changes,
confirming $\sigma=0.5$ is maximal.  Unlike the holomorphic eigenforms at
weights 12 and~20, no B-01 excess was observed---likely due to the wider
zero spacing ($\Delta t\approx 3.3$) of this Maass form.

\begin{table}[ht]
\centering
\caption{$\kappa\delta^{2}$ log-log slopes for the five lowest Maass zeros
  ($R=9.5337$, $\delta\in[0.005,0.1]$, 8 values).  All $R^{2}=1.000000$.}
\label{tab:maass_kd2}
\small
\begin{tabular}{llrl}
\toprule
Zero & $t$ & Slope & Rating \\
\midrule
$\rho_1$ & $3.264$ & $1.9999$ & $\bigstar\bigstar\bigstar$ \\
$\rho_2$ & $12.418$ & $2.0002$ & $\bigstar\bigstar\bigstar$ \\
$\rho_3$ & $17.577$ & $2.0004$ & $\bigstar\bigstar\bigstar$ \\
$\rho_4$ & $20.674$ & $2.0005$ & $\bigstar\bigstar\bigstar$ \\
$\rho_5$ & $23.553$ & $2.0007$ & $\bigstar\bigstar\bigstar$ \\
\midrule
Mean & & $2.0003\pm 0.0003$ & \\
\bottomrule
\end{tabular}
\end{table}

The slope $2.0003\pm 0.0003$ is the most precise P3 measurement across
all $L$-functions tested in Papers~1--3.  The exceptional precision stems
from the wide zero spacing: the smallest gap is $\rho_4$--$\rho_5$ at
$\Delta t=2.88$, much larger than typical holomorphic gaps at comparable
heights.  Consequently, proximate-zero interference---the dominant source
of error for weight-$16$ and weight-$20$ forms---is negligible.

\subsection{Even Maass Cusp Form ($R=13.7798$, Result~\#201)}
\label{ssec:maass_even}

Result~\#200 tested the first \emph{odd} Maass cusp form.  To verify
independence from the symmetry class and spectral parameter, we turn to
the first \emph{even} Maass cusp form on
$\mathrm{SL}_2(\mathbb{Z})\backslash\mathbb{H}$, with spectral parameter
$R=13.7797513518907389\ldots$ and $\varepsilon=+1$ (even symmetry class).
The completed $L$-function is:
\[
  \Lambda(s,u_R^{+}) = \pi^{-s}\,
  \Gamma\!\Bigl(\frac{s+iR}{2}\Bigr)
  \Gamma\!\Bigl(\frac{s-iR}{2}\Bigr)\,L(s,u_R^{+}), \qquad
  \Lambda(s)=\Lambda(1-s), \quad \varepsilon=+1.
\]

\medskip
\noindent\textbf{Result~\#201} (\textit{Even Maass cusp form, $R=13.7798$}):
\begin{itemize}
  \item P1 (FE): $|\Lambda(s)-\Lambda(1-s)|/|\Lambda(s)|=0.00$ at five test
        points (\texttt{dps}$=80$).
  \item P3 ($\kappa\delta^{2}$): Five-zero log-log slope
        $\mathbf{1.9999\pm 0.0008}$, $R^{2}=1.000000$ for all five zeros.
  \item P4 (monodromy): $5/5=2.0000\pi$.
  \item P5 ($\sigma$-uniqueness): FAIL (B-01); $\sigma=0.5\to 11$,
        $\sigma=0.1,0.2\to 12$.  Consistent with all $\mathrm{GL}(2)$ results.
\end{itemize}
Overall: \textbf{4/5 PASS} ($\bigstar\bigstar\bigstar$).
Combined with Result~\#200, the slope~$2$ persists across both even and
odd symmetry classes and two distinct spectral parameters ($R=9.53$,
$13.78$), confirming that the $\kappa\delta^{2}$ universality is
independent of the Maass spectral data.


\subsection{Degree Extension: $\mathrm{GL}(3)$ and $\mathrm{GL}(4)$}
\label{ssec:degree_ext}

All preceding $\kappa\delta^{2}$ measurements concern degree-$2$
$L$-functions (automorphic forms on $\mathrm{GL}(2)$).  A central
question is whether the slope~$2$ persists for higher-degree
$L$-functions in the Selberg class.  We test degree~$3$ and degree~$4$
via symmetric-power $L$-functions.

\paragraph{$\mathrm{GL}(3)$: $\mathrm{sym}^{2}(E, \text{11a1})$
(Result~\#202).}
\label{par:gl3sym2}
The symmetric-square $L$-function of the elliptic curve 11a1 is a
degree-$3$, conductor $N=121$ automorphic $L$-function with
$\Lambda(s)=\Lambda(3-s)$ and $\sigma_{\mathrm{crit}}=3/2$.
Five zeros ($t\in[3.9,12.1]$, spacing $>1.0$) are selected from
the $\approx 30$ zeros found by PARI's \texttt{lfunzeros}.

\medskip
\noindent\textbf{Result~\#202}: $\kappa\delta^{2}$ log-log slope
$\mathbf{2.0000\pm 0.0001}$, $R^{2}=1.000000$ for all five zeros.
Monodromy $5/5=2.0000\pi$.  $\sigma$-uniqueness: $5/5$ PASS
(centre $\sigma=1.5$ is maximal---unlike $\mathrm{GL}(2)$, no B-01
excess at degree~$3$).  FE: $|\Lambda(s)-\Lambda(3-s)|/|\Lambda(s)|=0.00$.
Overall: \textbf{5/5 PASS} ($\bigstar\bigstar\bigstar$).

The $\pm 0.0001$ error bar matches the best $\mathrm{GL}(2)$ measurement
(\#200, Maass odd).  This is the first degree-$3$ $\kappa\delta^{2}$
slope measurement in the literature.

\paragraph{$\mathrm{GL}(4)$: $\mathrm{sym}^{3}(\Delta)$
(Result~\#203).}
\label{par:gl4sym3}
The symmetric-cube $L$-function of the Ramanujan $\Delta$-function is a
degree-$4$, conductor $N=1$, motivic weight~$33$ automorphic $L$-function
with $\Lambda(s)=\Lambda(1-s)$ and $\sigma_{\mathrm{crit}}=1/2$ (PARI
normalisation with $\gamma_{V}=[5.5,6.5,16.5,17.5]$).  L(1/2)$\approx 0$
(root number $\varepsilon=-1$).  Five zeros ($t\in[4.2,12.1]$, spacing
$>1.0$) are found by manual $|L(1/2+it)|$ minimum search (golden-section
refinement, since \texttt{lfunzeros} is unreliable for large gamma shifts).

\medskip
\noindent\textbf{Result~\#203}: $\kappa\delta^{2}$ log-log slope
$\mathbf{2.0000\pm 0.0001}$, $R^{2}=1.000000$ for all five zeros.
Monodromy $5/5=2.0000\pi$.  $\sigma$-uniqueness: FAIL (centre count~$=24$,
max count~$=30$ at $\sigma=0.40$), consistent with the degree-${\geq}2$
structural FAIL pattern (B-05).  FE: \texttt{lfuncheckfeq}~$=-51$ digits.
Overall: \textbf{4/5 PASS} ($\bigstar\bigstar\bigstar$).

The individual zero slopes are $1.9999$, $2.0000$, $2.0001$, $1.9999$,
$2.0000$---each indistinguishable from the theoretical value~$2$.
Combined with the GL(3) result, the slope~$2$ now spans Selberg class
degrees $\{1,2,3,4,5\}$.

\paragraph{$\mathrm{GL}(4)$: $\mathrm{sym}^{3}(\text{37a1})$
(Result~\#204).}
\label{par:gl4sym3_37a1}
To address the potential bias toward the 11a1/\,$\Delta$ family
(sym$^{2}$(11a1) and sym$^{3}(\Delta)$ are both 11a1-derived), we
verify the same slope on a fully independent elliptic curve.
$L(s,\mathrm{sym}^3(\text{37a1}))$ has conductor $N=50653=37^3$,
degree~$4$, arithmetic normalisation $\sigma_{\mathrm{crit}}=k/2=2$,
and $\gamma_V=[-1,0,0,1]$ (PARI \texttt{lfunsympow(37a1,3)}).
FE: \texttt{lfuncheckfeq}~$=-137$ digits.

\medskip
\noindent\textbf{Result~\#204}: $\kappa\delta^{2}$ log-log slope
$\mathbf{1.9999\pm 0.0005}$, $R^{2}=1.000000$ for all five zeros
($t\in[0.37, 5.55]$, spacing $>1.0$).
Monodromy $5/5=2.0000\pi$.  $\sigma$-uniqueness: FAIL (centre=40,
max=70), consistent with degree-${\geq}2$ structural pattern (B-05).
Overall: \textbf{4/5 PASS} ($\bigstar\bigstar\bigstar$).

Individual slopes: $2.0002$, $2.0002$, $2.0000$, $2.0002$, $1.9988$.
This confirms that the slope~$2$ universality at degree~$4$ is
independent of the underlying elliptic curve, conductor, or motivic
weight ($w=33$ for $\mathrm{sym}^3(\Delta)$ vs.\ $w=3$ for
$\mathrm{sym}^3(\text{37a1})$).

\textcolor{ForestGreen}{[result \#204, confirmed 2026-04-21;
sym$^3$(37a1): $N{=}50653$, $k{=}4$, $\gamma_V{=}[-1,0,0,1]$;
FE${=}-137$ digits; slope${=}1.9999{\pm}0.0005$; mono $5/5{=}2\pi$;
$\sigma$-uniq FAIL (structural); 11-row table; 11a1-bias resolved]}


\paragraph{$\mathrm{GL}(5)$: $\mathrm{sym}^{4}(\text{11a1})$
(Result~\#205).}
\label{par:gl5sym4}
To extend the degree coverage to~$5$, we test the symmetric fourth
power $L(s,\mathrm{sym}^4(\text{11a1}))$: conductor $N=14641=11^4$,
degree~$5$, arithmetic normalisation $\sigma_{\mathrm{crit}}=k/2=2.5$,
$\gamma_V=[-2,-1,0,0,1]$ (PARI \texttt{lfunsympow(11a1,4)}).
FE: \texttt{lfuncheckfeq}~$=-116$ digits.

\medskip
\noindent\textbf{Result~\#205}: $\kappa\delta^{2}$ log-log slope
$\mathbf{1.9999\pm 0.0004}$, $R^{2}=1.000000$ for all five zeros
($t\in[1.49, 7.35]$, spacing $>1.0$).
Monodromy $5/5=2.0000\pi$.  $\sigma$-uniqueness: FAIL (centre=36,
max=68), consistent with degree-${\geq}2$ structural pattern (B-05).
Overall: \textbf{4/5 PASS} ($\bigstar\bigstar\bigstar$).

Individual slopes: $2.0000$, $1.9996$, $1.9995$, $2.0000$, $2.0006$.
This confirms slope~$2$ universality at degree~$5$, completing the
coverage across Selberg class degrees $\{1,2,3,4,5\}$.
Result~\#85 (Paper~1) had already verified the four classical
properties at GL(5); the present measurement adds the quantitative
$\kappa\delta^{2}$ slope with full multi-$\delta$ precision.
Degree~$5$ remains limited to one $L$-function due to computational
constraints (sym$^4$(37a1) has $N=37^4=1{,}874{,}161$).

\textcolor{ForestGreen}{[result \#205, confirmed 2026-04-21;
sym$^4$(11a1): $N{=}14641$, $k{=}5$, $\gamma_V{=}[-2,-1,0,0,1]$;
FE${=}-116$ digits; slope${=}1.9999{\pm}0.0004$; mono $5/5{=}2\pi$;
$\sigma$-uniq FAIL (structural); 12-row table; degree 1--5 complete]}


\subsection{Weight, Spectral, and Degree Invariance of $\kappa\delta^{2}$}
\label{ssec:weightinv}

Table~\ref{tab:weightuniv} collects the $\kappa\delta^{2}$ log-log slopes
(slope of $\log(\kappa\delta^{2}-1)$ vs.\ $\log\delta$) across all
weights, spectral parameters, and Selberg class degrees tested to date.

\begin{table}[ht]
\centering
\caption{Weight, spectral, degree, and family invariance of the $\kappa\delta^{2}$
  log-log slope.  Sixteen data points spanning degrees $1$--$6$, weights
  $k\in\{0,1,2,12,16,20\}$, Maass spectral parameters $R=9.53,13.78$,
  symmetric powers $\mathrm{sym}^{2}$--$\mathrm{sym}^{5}$,
  two Artin $L$-functions ($S_3$, $S_5$), one Rankin--Selberg
  convolution ($\text{11a1}\times\text{37a1}$), and one Dirichlet
  $L$-function ($\chi_{-7}$, conductor~$7$).
  All $\sigma$-direction measurements use the multi-$\delta$ log-log protocol.}
\label{tab:weightuniv}
\small
\begin{tabular}{llllrrl}
\toprule
Type & $L$-function & Degree & $\sigma_{\mathrm{crit}}$ & Slope & $\pm\sigma$ & Rating \\
\midrule
$k{=}0$  & Riemann $\zeta$ & $\mathrm{GL}(1)$ & $\tfrac{1}{2}$ & $\approx 2.0$ & --- & $\bigstar\bigstar\bigstar$ \\
Dirichlet & $L(s,\chi_{-7})$ (\#213) & $\mathrm{GL}(1)$ & $\tfrac{1}{2}$ & $\mathbf{2.0000}$ & $\mathbf{0.0014}$ & $\bigstar\bigstar\bigstar$ \\
$k{=}1$  & Artin $S_3$ (\#207) & $\mathrm{GL}(2)$ & $\tfrac{1}{2}$ & $\mathbf{2.0000}$ & $\mathbf{0.0000}$ & $\bigstar\bigstar\bigstar$ \\
$k{=}2$  & EC 11a1 & $\mathrm{GL}(2)$ & $1$ & $\approx 2.0$ & --- & $\bigstar\bigstar\bigstar$ \\
Maass odd & $u_{R}^{-}$ ($R{=}9.53$) & $\mathrm{GL}(2)$ & $\tfrac{1}{2}$ & $\mathbf{2.0003}$ & $\mathbf{0.0003}$ & $\bigstar\bigstar\bigstar$ \\
Maass even & $u_{R}^{+}$ ($R{=}13.78$) & $\mathrm{GL}(2)$ & $\tfrac{1}{2}$ & $1.9999$ & $0.0008$ & $\bigstar\bigstar\bigstar$ \\
$k{=}12$ & Ramanujan $\Delta$ & $\mathrm{GL}(2)$ & $6$ & $2.0008$ & $0.0006$ & $\bigstar\bigstar\bigstar$ \\
$k{=}16$ & $\Delta\cdot E_4$ & $\mathrm{GL}(2)$ & $8$ & $1.9989$ & $0.0035$ & $\bigstar\bigstar\bigstar$ \\
$k{=}20$ & $\Delta\cdot E_8$ & $\mathrm{GL}(2)$ & $10$ & $1.9984$ & $0.0055$ & $\bigstar\bigstar\bigstar$ \\
$\mathrm{sym}^{2}$ & $\mathrm{sym}^{2}(\text{11a1})$ & $\mathrm{GL}(3)$ & $\tfrac{3}{2}$ & $\mathbf{2.0000}$ & $\mathbf{0.0001}$ & $\bigstar\bigstar\bigstar$ \\
$\mathrm{sym}^{3}$ & $\mathrm{sym}^{3}(\Delta)$ & $\mathrm{GL}(4)$ & $\tfrac{1}{2}$ & $\mathbf{2.0000}$ & $\mathbf{0.0001}$ & $\bigstar\bigstar\bigstar$ \\
$\mathrm{sym}^{3}$ & $\mathrm{sym}^{3}(\text{37a1})$ & $\mathrm{GL}(4)$ & $2$ & $\mathbf{1.9999}$ & $\mathbf{0.0005}$ & $\bigstar\bigstar\bigstar$ \\
Artin $S_5$ & $\rho_4$ (\#209) & $\mathrm{GL}(4)$ & $\tfrac{1}{2}$ & $\mathbf{1.9999}$ & $\mathbf{0.0003}$ & $\bigstar\bigstar\bigstar$ \\
R-S & $\text{11a1}\times\text{37a1}$ (\#210) & $\mathrm{GL}(4)$ & $1$ & $\mathbf{1.9960}$ & $\mathbf{0.0072}$ & $\bigstar\bigstar\bigstar$ \\
$\mathrm{sym}^{4}$ & $\mathrm{sym}^{4}(\text{11a1})$ & $\mathrm{GL}(5)$ & $\tfrac{5}{2}$ & $\mathbf{1.9999}$ & $\mathbf{0.0004}$ & $\bigstar\bigstar\bigstar$ \\
$\mathrm{sym}^{5}$ & $\mathrm{sym}^{5}(\text{11a1})$ & $\mathrm{GL}(6)$ & $3$ & $\mathbf{1.9980}$ & $\mathbf{0.0036}$ & $\bigstar\bigstar\bigstar$ \\
\bottomrule
\end{tabular}
\end{table}

For holomorphic forms at $k=12$, $16$, and $20$ the slopes agree: the
three-weight mean is $1.9994\pm 0.0010$, and the maximum pairwise deviation
is $0.0024$.  The two Maass forms ($R=9.5337$ odd, $R=13.7798$ even) add
qualitatively different data points---non-holomorphic, weight~$0$, with
spectral gamma factors---yet their slopes ($2.0003\pm 0.0003$ and
$1.9999\pm 0.0008$) are among the most precise measurements.

The degree extension via symmetric powers is equally striking:
$\mathrm{sym}^{2}(\text{11a1})$ at degree~$3$ and $\mathrm{sym}^{3}(\Delta)$
at degree~$4$ both yield slope $2.0000\pm 0.0001$, matching the theoretical
value to four decimal places.  The Artin $S_5$ standard representation
$\rho_4$ (Result~\#209: irreducible degree-4, conductor~$2869$,
slope~$1.9999\pm 0.0003$) confirms universality beyond the
symmetric-power chain.  The Rankin--Selberg convolution
$L(s,\text{11a1}\times\text{37a1})$ (Result~\#210: degree~4, $N=407$,
slope~$1.9960\pm 0.0072$) provides the first tensor-product verification,
establishing a fourth independent degree-4 construction.
Result~\#213 adds four quadratic Dirichlet $L$-functions
$L(s,\chi_D)$ with $D\in\{-3,-4,5,-7\}$ (conductors $3$--$7$),
yielding a combined slope of $2.0002\pm 0.0001$---the most precise
measurement in the table.  This confirms that even the simplest
non-trivial $L$-functions (degree~$1$, $N>1$) satisfy the universal law,
and that conductor alone does not perturb the slope.
Combined with the $k=0$ (Riemann $\zeta$)
and $k=2$ (EC 11a1) results from Paper~1, the slope equals
$2.0\pm 0.008$ across sixteen data points spanning Selberg class degrees
$\{1,2,3,4,5,6\}$, holomorphic weights $k\in\{0,1,2,12,16,20\}$, two
non-holomorphic Maass forms, four Dirichlet $L$-functions,
four independent GL(4) $L$-functions
(two symmetric powers, one Artin, and one Rankin--Selberg),
one GL(5) $L$-function, and one GL(6) $L$-function.

\begin{theorem}[Slope universality]
\label{thm:slopeuniv}
  Let $L(s,\pi)$ be an automorphic $L$-function with completed function
  $\Lambda(s,\pi)$ and contragredient $\Lambda(s,\tilde\pi)$ satisfying:
  \begin{enumerate}
    \item[\textup{(FE)}] $\Lambda(s,\pi) = \varepsilon\,\Lambda(k-s,\tilde\pi)$
          for some $k\in\mathbb{R}$, $|\varepsilon|=1$; and
    \item[\textup{(SR)}] $\Lambda(\bar{s},\tilde\pi) = \overline{\Lambda(s,\pi)}$
          \textup{(}Schwarz reflection\textup{)}.
  \end{enumerate}
  Let $\rho=\tfrac{k}{2}+i\gamma$ be a simple zero on the critical line.
  Write the Laurent expansion
  $\Lambda'/\Lambda(s) = (s-\rho)^{-1} + c_0 + c_1(s-\rho) + \cdots$.
  Then:
  \begin{enumerate}
    \item[\textup{(i)}] $\operatorname{Re}(c_0)=0$ \textup{(}$c_0$ is purely imaginary\textup{)};
    \item[\textup{(ii)}] $\operatorname{Im}(c_1)=0$ \textup{(}$c_1$ is real\textup{)};
    \item[\textup{(iii)}] $\kappa(\rho{+}\delta)\,\delta^{2}
          = 1 + A\,\delta^{2} + O(\delta^{3})$
          with $A = \operatorname{Im}(c_0)^{2} + 2c_1$.
  \end{enumerate}
  In particular, $\log(\kappa\delta^{2}-1) = \log A + 2\log\delta + O(\delta)$,
  so the log-log slope is exactly~$2$.
\end{theorem}

\begin{proof}
  Write Laurent coefficients $c_j$ for $\Lambda'/\Lambda(s,\pi)$ at $\rho$
  and $c_j'$ for $\Lambda'/\Lambda(s,\tilde\pi)$ at $\bar\rho=k-\rho$.
  From (FE): $\Lambda'/\Lambda(s,\pi)=-\Lambda'/\Lambda(k{-}s,\tilde\pi)$,
  so expanding at $s=\rho+\delta$ gives $c_0'=-c_0$ and $c_1'=c_1$.
  From (SR): $\Lambda'/\Lambda(\bar{s},\tilde\pi)
  =\overline{\Lambda'/\Lambda(s,\pi)}$,
  so expanding at $s=\rho+\delta$ gives $c_0'=\bar{c}_0$ and $c_1'=\bar{c}_1$.
  Combining: $c_0=-\bar{c}_0$ (so $\operatorname{Re}(c_0)=0$)
  and $c_1=\bar{c}_1$ (so $\operatorname{Im}(c_1)=0$).
  Then $\kappa\delta^{2}=|1+c_0\delta+c_1\delta^2|^2
  =1+A\delta^2+O(\delta^3)$,
  since $\operatorname{Re}(c_0)=0$ kills the $\delta^1$ term.
\end{proof}

\begin{proposition}[Laurent coefficient parity---universal]
\label{prop:laurentparity}
The FE/SR argument yields $c_n = (-1)^{n+1}\bar{c}_n$ for
\emph{every} Laurent coefficient $c_n$, $n\geq 0$,
and for \emph{every} $L$-function satisfying
$\Lambda(s,\pi)=\varepsilon\,\Lambda(1{-}s,\tilde{\pi})$ and
$\Lambda(\bar{s},\pi)=\overline{\Lambda(s,\tilde{\pi})}$.
Self-duality ($\pi=\tilde{\pi}$) is \textbf{not required}:
the FE relates $c_n(\pi,\rho)$ to $c_n(\tilde{\pi},\bar{\rho})$,
and SR relates $c_n(\pi,\rho)$ to $\bar{c}_n(\tilde{\pi},\bar{\rho})$;
combining gives $c_n=(-1)^{n+1}\bar{c}_n$ for any~$\pi$.
Consequently, $c_{2k}$ is purely imaginary and $c_{2k+1}$ is real
for all $k\geq 0$.
Numerical verification at orders $n=0$--$3$:
$11$~self-dual zeros ($\zeta$, 11a1, $\chi_{-3}$; Result~\#C-240)
and $6$~non-self-dual zeros ($\chi\bmod 5$, $\chi\bmod 7$;
Result~\#C-241) all satisfy $|\operatorname{Re}(c_{2k})|/|c_{2k}|\leq
3\times 10^{-7}$ and $\operatorname{Im}(c_{2k+1})=0$ to machine
precision ($200$-digit arithmetic).
\textbf{Sharpness.}
The hypothesis $\sigma=\tfrac{1}{2}$ (i.e.\ $1{-}\rho=\bar{\rho}$)
is essential: the Davenport--Heilbronn function satisfies both FE and
SR, yet its four known off-critical zeros
($\sigma\neq\tfrac{1}{2}$) show
$|\!\operatorname{Re}(c_0)|/|\!\operatorname{Im}(c_0)|\sim 2$--$7$,
while five on-critical zeros satisfy $\sim 10^{-12}$
(Result~\#C-242).
Thus parity characterises zeros on the critical line.
\end{proposition}

\begin{proof}
The functional equation maps $c_n(\pi,\rho)$ to $\varepsilon(-1)^n c_n(\tilde{\pi},\bar{\rho})$.
The conjugate-centric symmetry maps $c_n(\pi,\rho)$ to $\bar{c}_n(\tilde{\pi},\bar{\rho})$.
Combining: $c_n = (-1)^{n+1}\bar{c}_n$.
\end{proof}

\begin{theorem}[Off-critical parity protection]
\label{thm:parityprotection}
Under the hypotheses of Proposition~\textup{\ref{prop:laurentparity}},
let $\rho=(\tfrac{1}{2}+\delta)+it_0$ with $\delta\neq 0$ be a simple
off-critical zero, and let $c_n$ denote the Laurent coefficients of
$\Lambda'/\Lambda$ at~$\rho$.
Then:
\begin{enumerate}[\upshape(a)]
\item \textbf{Odd-order protection.}
      $r(c_{2k+1})\coloneqq |\!\operatorname{Im}(c_{2k+1})|/|\!\operatorname{Re}(c_{2k+1})|
       = O(|\delta|^{2k+3})$ as $\delta\to 0$.
\item \textbf{Even-order amplification.}
      $r(c_{2k})\coloneqq |\!\operatorname{Re}(c_{2k})|/|\!\operatorname{Im}(c_{2k})|
       = O(|\delta|^{-(2k+1)})$ as $\delta\to 0$.
\end{enumerate}
In particular, the parity violation for odd-order coefficients decays
as a power law in the off-critical displacement $|\delta|$,
with the exponent growing linearly in the order~$k$.
\end{theorem}

\begin{proof}
By the FE and conjugate-centric symmetry,
the zero~$\rho$ has a partner
$1{-}\bar{\rho}=(\tfrac{1}{2}-\delta)+it_0$.
In the Hadamard product representation, the partner zero at distance
$2\delta$ contributes $(-1)^n/(2\delta)^{n+1}$ (a real quantity) to $c_n$.
For odd~$n$, this dominates $\operatorname{Re}(c_n)$ (the parity-preserving
component), giving $|\!\operatorname{Re}(c_{2k+1})|\asymp (2\delta)^{-(2k+2)}$.
The parity-violating component $\operatorname{Im}(c_{2k+1})$ arises from
the first-order $\delta$-perturbation of on-critical zero contributions
in the Hadamard sum; its leading term is $O(\delta)$.
The ratio gives $r(c_{2k+1})=O(\delta/(2\delta)^{-(2k+2)})=O(\delta^{2k+3})$.
For even~$n$, the partner contributes $O((2\delta)^{-(2k+1)})$
to $\operatorname{Re}(c_{2k})$ (the parity-\emph{violating} component),
while $|\!\operatorname{Im}(c_{2k})|$ is dominated by
on-critical-zero contributions of order $O(1)$, so $r(c_{2k})$ grows.
\end{proof}

\noindent\textit{Numerical verification} (Result~\#C-243):
Laurent coefficients $c_0$--$c_7$ computed at $5$~on-critical and
$4$~off-critical zeros of the Davenport--Heilbronn function
($\mathrm{dps}=120$, $72$~values total).
The \emph{partner dominance ratio}
$|\!\operatorname{Re}(c_{2k+1})|/(2\delta)^{-(2k+2)}$
approaches~$1$ as $\delta\to 0$
($0.98$ for $c_1$ at $\delta=0.074$; $1.00$ for $c_3$).
The $\delta$-exponents $\beta_k\coloneqq \log r/\log|\delta|$
average $2.91$, $4.38$, $5.94$, $7.74$ for $k=0,1,2,3$
(predicted: $3,5,7,9$); the systematic deficit is attributed to
$t$-dependent prefactors confounded with~$\delta$.
Even/odd asymmetry ratio: $2.4\times 10^7$.

\begin{corollary}[Amplitude formula]
\label{cor:amplitude}
  The sub-leading coefficient $A=\operatorname{Im}(c_0)^2+2c_1$ decomposes
  into a \emph{local density} contribution $\operatorname{Im}(c_0)^2$
  (related to nearby zeros via $c_0=\sum_{\rho'\neq\rho}(\rho-\rho')^{-1}$)
  and a \emph{gamma-factor} contribution $2c_1$ (dominated by the digamma
  terms in $\Lambda'/\Lambda$).
\end{corollary}

\noindent\textit{Numerical verification} (Result~\#214):
$35$ zeros across seven $L$-functions---four self-dual
($\zeta$, $L(s,\chi_{-3})$, $L(s,\chi_{-7})$, $L(s,\chi_{-11})$)
and three non-self-dual
($\chi\bmod 5$, $\chi\bmod 7$, $\chi\bmod 13$
of orders $4$, $6$, $12$)---yield
$|\operatorname{Re}(c_0)|=0$ to machine precision ($150$-digit arithmetic),
mean slope $2.0004\pm 0.0023$, and mean $A$-prediction error $0.6\%$
(max $1.7\%$; Result~\#216).
The amplitude formula is verified at higher degrees (Result~\#217):
$25$~zeros across five $L$-functions spanning GL(1)--GL(4)
($\zeta$, two elliptic curves, $\mathrm{sym}^{2}$, $\mathrm{sym}^{3}$)
yield $A$-prediction errors of mean $1.3\%$ (max $4.0\%$ at degree~$4$).
An independent cross-validation (Result~\#218) adds $20$~zeros from
four \emph{distinct} families---two Dirichlet $L$-functions ($\chi_{-3}$,
$\chi_{-7}$), the Ramanujan $\Delta$ $L$-function (weight~$12$, non-EC),
and $\mathrm{sym}^{2}(\text{37a1})$ (rank-$1$ curve, independent of~11a1)---all
yielding $20/20$ PASS with mean error $0.5\%$ (max $2.0\%$).
Together, $45$~zeros across nine $L$-functions confirm
Corollary~\ref{cor:amplitude} across degrees $1$--$4$,
ruling out symmetric-power chain bias.
The theorem subsumes the earlier Theorem~\ref{thm:rankinv} (which
established $\kappa\delta^2=1+O(\delta^2)$ without the explicit $A$ formula)
and upgrades the empirical observation of slope~$2$ in Table~\ref{tab:weightuniv}
from a numerical regularity to a mathematical theorem.
Sixteen data points spanning Selberg class degrees $\{1,2,3,4,5,6\}$,
holomorphic weights $k\in\{0,1,2,12,16,20\}$, two Maass forms,
four Dirichlet $L$-functions, four independent GL(4) constructions
(symmetric powers, Artin, Rankin--Selberg), and one GL(6) $L$-function
all confirm the predicted slope~$2$ to within $\pm 0.008$,
establishing \emph{construction universality}.

\paragraph{Amplitude scaling: the sub-leading coefficient $A(t_{0})$.}
While the slope is a universal constant ($=2$), the amplitude
$A(t_{0}) \;\defeq\; \kappa(\sigma_{\mathrm{crit}}+\delta,\,t_{0}) - 1/\delta^{2}$
at small $\delta$ varies across zeros and $L$-functions.
Extracting $A$ at $\delta=0.01$ from $52$ zeros across degrees $1$--$6$
(Result~\#212) reveals a two-variable dependence:
\begin{equation}\label{eq:Ascaling}
  A(d,N) \;\approx\; \alpha\,d^{2} \;+\; \beta\,\log N \;+\; \gamma,
  \qquad \alpha\approx 0.39,\;\beta\approx 2.19,\;\gamma\approx 0.10,
\end{equation}
with in-sample $R^{2}=0.983$.  The degree-squared term reflects the
Hadamard product decomposition
$A = B(t_{0})^{2} + 2H_{1}(t_{0})$,
where $H_{1}=\sum_{\rho'\neq\rho}(\mathrm{Im}\,\rho'-t_{0})^{-2}$
grows with zero density (hence with $d$), and the $\log N$ term
arises from the Stirling approximation of the gamma factors in the
completed $L$-function.

\emph{Out-of-sample validation.}  For the Artin $S_5$ $L$-function
(Result~\#209: $d=4$, $N=2869$), equation~\eqref{eq:Ascaling} predicts
$A\approx 23.9$; the measured mean over five zeros is $24.9$ (error~$4\%$).
For $L(s,\text{11a1}\times\text{37a1})$ (Result~\#210: $d=4$, $N=407$),
the prediction is $19.6$ versus measured~$30.7$ (high variance due to a
close zero pair).  The qualitative trend---$A$ growing with both $d^{2}$
and $\log N$---is robust; the precise coefficients require further data
for stabilization.
\textcolor{ForestGreen}{[result \#212, confirmed 2026-04-21;
6 degrees, 52 zeros; out-of-sample Artin $S_5$ validation $4\%$ error;
Observation: $A(d,N)\propto d^{2}+\log N$ (tentative coefficients)]}

\emph{Intrinsic variability of $A(t_0)$.}  Direct computation of $H_1$
via the Hadamard sum for $269$ zeros of $\zeta(s)$ (Result~\#215) reveals
a coefficient of variation $\mathrm{CV}(H_1)\approx 42\%$, arising from
the GUE-type fluctuations in local zero spacing.  The normalised ratio
$H_1\cdot\Delta(t_0)^2$ (where $\Delta=2\pi/\log(t_0/2\pi)$ is the mean
spacing) has a stable mean of $4.58\pm 1.95$ across $t_0\in[50,500]$,
but the point-to-point variability precludes a precise universal constant.
The $B(t_0)^2$ term (from the Hilbert-type sum
$\sum_{\rho'\neq\rho}(\mathrm{Im}\,\rho'-t_0)^{-1}$) contributes
${\approx}24\%$ of $A$, fluctuating with mean zero and standard deviation
${\approx}1.0$.  This intrinsic $O(1)$ variability explains why the
regression~\eqref{eq:Ascaling} achieves $R^2=0.98$ in-sample but remains
``tentative'' for point predictions.
\textcolor{ForestGreen}{[result \#215, confirmed 2026-04-21;
269 zeros of $\zeta$; $H_1=\sum 1/(\gamma_0-\gamma_n)^2$ verified
vs PARI to $0.7\%$; CV$=42\%$; $2H_1$ contributes $76\%$ of $A$]}
\textcolor{ForestGreen}{[result \#217, confirmed 2026-04-21;
$A$ formula verified degrees $1$--$4$; $25$ zeros, $5$ $L$-functions;
mean error $1.3\%$, max $4.0\%$]}


% ============================================================================
\section{Discussion}
\label{sec:disc}
% ============================================================================

\paragraph{Universality.}
The four $\xi$-bundle properties pass across an extended family of
$L$-functions: elliptic-curve $L$-functions (rank 0--3, conductor
11--10099, 8 curves, $\varepsilon=\pm 1$), Rankin--Selberg
$\mathrm{GL}(4)$ (2 pairs), Dedekind $\zeta$-functions (7 imaginary
quadratic fields), and Epstein $\zeta$-functions including 3 non-Euler
forms (class number 2--3).  Combined with Paper~1's $\mathrm{GL}(1)$--
$\mathrm{GL}(5)$ results (81 numerical experiments), the framework now
covers the major algebraic families of $L$-functions.

\paragraph{Weight, spectral, and degree universality.}
Results~\#125--\#127, \#200--\#201, and \#202--\#203 establish that the
$\kappa\delta^{2}$ log-log slope is invariant under the modular weight,
holomorphicity, Maass symmetry class, and Selberg class degree.
For holomorphic eigenforms the three-weight mean is $1.9994\pm 0.0010$
($k=12,16,20$); the two Maass cusp forms yield $2.0003\pm 0.0003$
($R=9.53$, odd) and $1.9999\pm 0.0008$ ($R=13.78$, even).  The degree
extension via symmetric powers produces the most precise measurements:
$\mathrm{sym}^{2}(\text{11a1})$ at $\mathrm{GL}(3)$ and
$\mathrm{sym}^{3}(\Delta)$ at $\mathrm{GL}(4)$ both give slope
$2.0000\pm 0.0001$.  Combined with the $k=0,2$ results from Paper~1,
the Artin result ($k=1$, Paper~3), and the Rankin--Selberg tensor product
(Result~\#210, slope $1.9960\pm 0.0072$), the slope equals
$2.0\pm 0.008$ across sixteen data points spanning Selberg class degrees
$\{1,2,3,4,5,6\}$ and six construction families.
Theorem~\ref{thm:slopeuniv} proves this analytically: the functional
equation and reality condition force $\operatorname{Re}(c_0)=0$ in the
Laurent expansion, killing the $\delta^1$ term and making slope~$2$ a
mathematical necessity for all automorphic $L$-functions.

\paragraph{Euler product dispensability.}
Theorem~\ref{thm:feonly} and Result~\#114 establish that the Euler
product is not required for the $\xi$-bundle properties.  This aligns
with the interpretation that the $\xi$-bundle geometry is a consequence
of the functional equation structure alone.

\paragraph{The discrimination theorem.}
Theorem~\ref{thm:disc} provides the first \emph{analytic} statement
connecting the $O(\delta)$ correction to $\kappa\delta^{2}$ with the
location of a zero relative to the critical line.  The numerical
verification (Results~\#115--\#117) confirms this connection for the
only known explicit test case (DH function).  For the Riemann
$\zeta$-function, no off-critical zeros are known; Theorem~\ref{thm:disc}
implies that if such a zero existed, it would produce a detectable
$O(\delta)$ signal.

\paragraph{Limitations.}
The asymptotic $c_{1}\cdot d \to 1$ is validated with $N=4$ data points.
The systematic $\delta$-sweep (Result~\#116) is limited to $N=3$
off-critical zeros (OFF\#4 had insufficient precision at $\delta<0.1$).
No new off-critical zeros were discovered in the range $t\in[200,400]$
(22 candidate points checked, 0 confirmed), consistent with the
known scarcity of DH off-critical zeros.

\paragraph{Future directions.}
\begin{enumerate}
  \item \textit{Hilbert modular forms}: extending the spectral universality
        to number fields of degree $>1$.
  \item \textit{Higher-degree extension}: degree~$7+$ via
        $\mathrm{sym}^{6}$ or other automorphic constructions.
  \item \textit{Selberg class generalisation}: Theorem~\ref{thm:disc}
        uses only FE; can it be sharpened to give $c_{1}$ in terms of
        the Selberg class degree?
  \item \textit{$\sigma$-localisation}: extending the discrimination to
        non-horizontal offsets (complex $\delta$).
  \item \textit{Additional DH off-critical zeros}: extending the search
        to $t>400$ with higher precision.
\end{enumerate}


% ============================================================================
\section{Conclusion}
\label{sec:conclusion}
% ============================================================================

We have extended the $\xi$-bundle framework of Kim~(2026) in four
directions.  First, universality: eleven new numerical results
(\#107--\#117) confirm the four $\xi$-bundle properties across elliptic-
curve, Rankin--Selberg, Dedekind, and Epstein $L$-functions; in
particular, non-Euler Epstein $\zeta$-functions (Theorem~6) demonstrate
that the Euler product is not necessary.  Second, discrimination:
Theorem~7 proves analytically that the first Taylor coefficient
$c_{1}$ of $\kappa\delta^{2}$ vanishes if and only if the zero lies on
the critical line; numerical verification via $\delta$-sweep and direct
analytic computation confirms this for all four known DH off-critical
zeros.  Third, weight and spectral universality: Results~\#125--\#127
and \#200--\#201 establish that the slope equals $2.0\pm 0.001$
independently of weight ($k=12,16,20$), holomorphicity (two Maass forms,
$R=9.53$ and $13.78$), and symmetry class (odd and even).
Fourth, degree and family universality: Results~\#202--\#206 extend the slope
verification from $\mathrm{GL}(2)$ through $\mathrm{GL}(6)$
via symmetric powers; Result~\#209 adds an irreducible degree-4
Artin $L$-function (standard representation of~$S_5$, slope
$1.9999\pm 0.0003$); and Result~\#210 adds the Rankin--Selberg
tensor product $L(s,\text{11a1}\times\text{37a1})$ (slope
$1.9960\pm 0.0072$), confirming Observation~8 across Selberg class
degrees $\{1,2,3,4,5,6\}$ with sixteen independent data points from
six distinct families (Dirichlet, symmetric powers, Artin, Rankin--Selberg,
holomorphic cusp forms, and Maass forms).

Together, Theorems~5--7 and Observation~8 form a coherent picture:
on the critical line, the $\xi$-bundle curvature satisfies
$\kappa\delta^{2}=1+O(\delta^{2})$ universally across all $L$-function
families, modular weights, and spectral types tested; off the critical
line, a qualitatively different $O(\delta)$ correction appears,
analytically traced to the breaking of mirror symmetry by the functional
equation.


% ============================================================================
% APPENDIX
% ============================================================================
\appendix

\section{Summary Table: All Thirty-One Results}
\label{app:summary}

\begin{longtable}{lllp{5.5cm}}
\caption{Summary of Results~\#107--\#117, \#125--\#127, \#200--\#218, \#C-240--\#C-242 (34 results). Rating: \textbf{$\bigstar\bigstar\bigstar$} strong
  positive, \textbf{$\bigstar\bigstar$} positive, \textbf{C} conditional.}
\label{tab:summary}\\
\toprule
Result & Section & Rating & Key finding \\
\midrule
\endfirsthead
\multicolumn{4}{c}{\tablename~\thetable{} (continued)}\\
\toprule
Result & Section & Rating & Key finding \\
\midrule
\endhead
\bottomrule
\endfoot
\#107 & \S\ref{ssec:epsilon}      & \textbf{$\bigstar\bigstar$}\,C  & EC 8 curves; P4 fails for $\varepsilon=-1$ (Hardy $Z$ method); P3 passes all \\
\#108 & \S\ref{ssec:epsilon}      & \textbf{$\bigstar\bigstar\bigstar$}    & EC 8 curves, $\varepsilon$-corrected; all 4/4 PASS; $\kappa\delta^{2}=1.000$ \\
\#109 & \S\ref{ssec:thm5}         & \textbf{$\bigstar\bigstar$}     & $\kappa_{\text{near}}$ universality; CV $<0.13\%$ at $\delta=0.01$; rank/conductor independent \\
\#110 & \S\ref{ssec:mono}         & \textbf{$\bigstar\bigstar\bigstar$}    & 8 EC curves, 120 zeros; $|\Delta\!\arg/\pi|=1.000000$; contour $=2.000000\pi$ \\
\#111 & \S\ref{ssec:conductor}    & \textbf{$\bigstar\bigstar$}     & Conductor scaling $\Delta\kappa\approx 2\log N\times 10^{-3}$; $R^{2}=0.73$; effect $<2\%$ \\
\#112 & \S\ref{sec:gl4rs}         & \textbf{$\bigstar\bigstar$}     & R-S GL(4), 2 pairs; 4/4 PASS; $\kappa\delta^{2}=1.014$, $1.033$ \\
\#113 & \S\ref{sec:dedekind}      & \textbf{$\bigstar\bigstar\bigstar$}    & Dedekind $\zeta$, 7 fields (h=1,3); 4/4 PASS; all $\kappa\delta^{2}<1.005$ \\
\#114 & \S\ref{sec:epstein}       & \textbf{$\bigstar\bigstar\bigstar$}    & Epstein $\zeta$; 3 non-Euler (h=2,3); 4/4 PASS; FE-only sufficiency (Thm~6) \\
\#115 & \S\ref{ssec:dhtest}       & \textbf{$\bigstar\bigstar\bigstar$}    & DH on/off discrimination; $1.002\pm0.001$ vs.\ $1.210\pm0.099$; complete separation \\
\#116 & \S\ref{ssec:sweep}        & \textbf{$\bigstar\bigstar\bigstar$}    & $\delta$-sweep; slope $1.999\pm0.003$ (on) vs.\ $0.963\pm0.039$ (off); Thm~7 confirmed \\
\#117 & \S\ref{ssec:c1}           & \textbf{$\bigstar\bigstar$}     & $c_{1}=\mathrm{Re}(\Lambda''/\Lambda')$ verified ($<0.5\%$, $N=4$); $c_{1}\cdot d\to 1$ \\
\#125 & \S\ref{ssec:delta12}      & \textbf{$\bigstar\bigstar\bigstar$}    & Ramanujan $\Delta$ ($k=12$); 4/5 PASS; $\kappa\delta^{2}$ slope $2.0008\pm0.0006$; $\sigma$-uniq.\ FAIL (B-01) \\
\#126 & \S\ref{ssec:w16}          & \textbf{$\bigstar\bigstar\bigstar$}    & Weight-16 cusp form ($\Delta\cdot E_4$); 5/5 PASS; slope $1.9989\pm0.0035$; weight invariance (Obs.~8) \\
\#127 & \S\ref{ssec:w20}          & \textbf{$\bigstar\bigstar\bigstar$}    & Weight-20 cusp form ($\Delta\cdot E_8$); 4/5 PASS; slope $1.9984\pm0.0055$; $\sigma$-uniq.\ FAIL (B-01) \\
\#200 & \S\ref{ssec:maass}        & \textbf{$\bigstar\bigstar\bigstar$}    & Maass cusp form ($R{=}9.5337$, odd); 5/5 PASS; slope $\mathbf{2.0003\pm0.0003}$ \\
\#201 & \S\ref{ssec:maass_even}   & \textbf{$\bigstar\bigstar\bigstar$}    & Even Maass ($R{=}13.7798$); 4/5 PASS; slope $1.9999\pm0.0008$; symmetry-class independence \\
\#202 & \S\ref{ssec:degree_ext}   & \textbf{$\bigstar\bigstar\bigstar$}    & GL(3) $\mathrm{sym}^{2}(\text{11a1})$; 5/5 PASS; slope $\mathbf{2.0000\pm0.0001}$; first degree-3 measurement \\
\#203 & \S\ref{ssec:degree_ext}   & \textbf{$\bigstar\bigstar\bigstar$}    & GL(4) $\mathrm{sym}^{3}(\Delta)$; 4/5 PASS; slope $\mathbf{2.0000\pm0.0001}$; degree 1--4 universality \\
\#204 & \S\ref{ssec:degree_ext}   & \textbf{$\bigstar\bigstar\bigstar$}    & GL(4) $\mathrm{sym}^{3}(\text{37a1})$; 4/5 PASS; slope $\mathbf{1.9999\pm0.0005}$; 11a1-bias resolved \\
\#205 & \S\ref{ssec:degree_ext}   & \textbf{$\bigstar\bigstar\bigstar$}    & GL(5) $\mathrm{sym}^{4}(\text{11a1})$; 4/5 PASS; slope $\mathbf{1.9999\pm0.0004}$; degree 1--5 complete \\
\#206 & \S\ref{ssec:degree_ext}   & \textbf{$\bigstar\bigstar\bigstar$}    & GL(6) $\mathrm{sym}^{5}(\text{11a1})$; 4/5 PASS; slope $\mathbf{1.9980\pm0.0036}$; degree 1--6 complete \\
\#209 & \S\ref{ssec:degree_ext}   & \textbf{$\bigstar\bigstar\bigstar$}    & Artin $S_5$ $\rho_4$; 4/4 PASS; slope $\mathbf{1.9999\pm0.0003}$; independent degree-4 (Paper~3) \\
\#210 & \S\ref{sec:gl4rs}         & \textbf{$\bigstar\bigstar\bigstar$}    & R-S $\text{11a1}\times\text{37a1}$; eff.\ 4/4 PASS; slope $\mathbf{1.9960\pm0.0072}$; tensor-product construction \\
\#212 & \S\ref{ssec:weightinv}    & \textbf{$\bigstar\bigstar$}           & $A(d,N)$ scaling law; $R^{2}=0.983$; Artin $S_5$ out-of-sample $4\%$ error; tentative coefficients \\
\#213 & \S\ref{ssec:weightinv}    & \textbf{$\bigstar\bigstar\bigstar$}    & Dirichlet $L(s,\chi_D)$, $D\in\{-3,-4,5,-7\}$; slope $\mathbf{2.0002\pm0.0001}$; 6th family (conductor independence) \\
\#215 & \S\ref{ssec:weightinv}    & \textbf{$\bigstar\bigstar$}           & $A(t_0)$ Hadamard decomposition; $H_1$ verified $0.7\%$; CV$=42\%$ intrinsic; $2H_1=76\%$ of $A$ \\
\#216 & \S\ref{ssec:weightinv}    & \textbf{$\bigstar\bigstar\bigstar\bigstar$} & Non-self-dual ($\chi_5,\chi_7,\chi_{13}$); slope $\mathbf{2.0004\pm0.0024}$; FE+SR proof; Thm generalised \\
\#217 & \S\ref{ssec:weightinv}    & \textbf{$\bigstar\bigstar\bigstar\bigstar$} & $A$ formula degrees $1$--$4$; $25$ zeros, $5$ $L$-functions; mean error $1.3\%$; Cor.~\ref{cor:amplitude} verified \\
\#218 & \S\ref{ssec:weightinv}    & \textbf{$\bigstar\bigstar\bigstar\bigstar$} & $A$ formula cross-validation; $20/20$ PASS; $\Delta$, $\chi_{-3}$, $\chi_{-7}$, $\mathrm{sym}^{2}(\text{37a1})$; mean error $0.5\%$; chain bias ruled out \\
\#C-240 & \S\ref{ssec:weightinv} & \textbf{$\bigstar\bigstar\bigstar$} & Laurent parity $c_n=(-1)^{n{+}1}\bar{c}_n$; $11/11$ PASS (self-dual); $|\!\operatorname{Re}(c_2)|/|c_2|\leq 3{\times}10^{-7}$; Proposition~\ref{prop:laurentparity} \\
\#C-241 & \S\ref{ssec:weightinv} & \textbf{$\bigstar\bigstar\bigstar\bigstar$} & Parity \textbf{universal}: non-self-dual $\chi_5,\chi_7$ also $6/6$ PASS; self-dual not needed (FE+SR suffice); Proposition~\ref{prop:laurentparity} \\
\#C-242 & \S\ref{ssec:weightinv} & \textbf{$\bigstar\bigstar\bigstar\bigstar$} & Parity \textbf{sharpness}: DH on-critical $5/5$ PASS ($r(c_0){\sim}10^{-12}$), off-critical $4/4$ FAIL ($r(c_0){\sim}4.4$); $\sigma{=}\tfrac{1}{2}$ essential; Proposition~\ref{prop:laurentparity} \\
\#C-243 & \S\ref{ssec:weightinv} & \textbf{$\bigstar\bigstar\bigstar\bigstar$} & Parity \textbf{protection}: $r(c_{2k{+}1}){=}O(|\delta|^{2k{+}3})$; partner dominance $0.98$--$1.00$; even/odd ratio $2.4{\times}10^{7}$; Theorem~\ref{thm:parityprotection} \\
\end{longtable}


\section{Numerical Data: Off-Critical Zeros (DH)}
\label{app:dhzeros}

Table~\ref{tab:dhdata} collects the complete $\delta$-sweep data for the
three DH off-critical zeros with sufficient precision (OFF\#4 included
only at $\delta\geq 0.10$).

\begin{table}[ht]
\centering
\caption{$\kappa\delta^{2}$ values from Result~\#116, DH off-critical zeros.
  ON-critical ensemble mean (10 zeros) shown for comparison.}
\label{tab:dhdata}
\small
\begin{tabular}{lrrrrrrrr}
\toprule
Zero & $\delta{=}0.005$ & $0.01$ & $0.015$ & $0.02$ & $0.03$ & $0.05$ & $0.07$ & $0.10$ \\
\midrule
ON (mean) & 1.000048 & 1.000192 & 1.000432 & 1.000767 & 1.001726 & 1.004795 & 1.009399 & 1.019182 \\
OFF\#1 & 1.018828 & 1.037696 & 1.056606 & 1.075561 & 1.113605 & 1.190280 & 1.267813 & 1.385945 \\
OFF\#2 & 1.034427 & 1.068469 & 1.102132 & 1.135420 & 1.200904 & 1.327729 & 1.449538 & 1.624253 \\
OFF\#3 & 1.067034 & 1.132039 & 1.195056 & 1.256139 & 1.372769 & 1.585968 & 1.776429 & 2.028926 \\
OFF\#4 & --- & --- & --- & --- & --- & --- & --- & 1.506337 \\
\bottomrule
\end{tabular}
\end{table}


\section{Mirror-Pair Decomposition of $c_{1}$}
\label{app:mirror}

For an off-critical zero $\rho=\sigma_{0}+it_{0}$, the mirror partner
is $\tilde\rho=1-\sigma_{0}+it_{0}$.  The leading contribution of the
mirror pair to $A'(\rho)/A(\rho)$ is
\[
  \frac{1}{\rho-\tilde\rho} = \frac{1}{2\sigma_{0}-1} = \frac{1}{2d},
\]
giving $c_{1}^{\text{mirror}} = 2\,\mathrm{Re}(1/(2d)) = 1/d$.
Table~\ref{tab:mirror} shows the decomposition from Result~\#117.

\begin{table}[ht]
\centering
\caption{Mirror-pair vs.\ background decomposition of $c_{1}$ (Result~\#117).}
\label{tab:mirror}
\small
\begin{tabular}{lrrrrrr}
\toprule
Zero & $d$ & $c_{1}^{\text{fit}}$ & $c_{1}^{\text{mirror}}=1/d$ & $c_{1}^{\text{bg}}$ & $c_{1}^{\text{bg}}/c_{1}^{\text{mirror}}$ \\
\midrule
OFF\#1 & 0.3085 & 3.760 & 3.241 & 2.139 & 1.320 \\
OFF\#2 & 0.1508 & 6.916 & 6.630 & 3.601 & 1.086 \\
OFF\#3 & 0.0744 & 13.543 & 13.449 & 6.818 & 1.014 \\
OFF\#4 & 0.2243 & 5.026 & 4.459 & 2.796 & 1.254 \\
\bottomrule
\end{tabular}
\end{table}

As $d\to 0$, the background-to-mirror ratio approaches $1$, indicating
that the two contributions are of equal magnitude in the limit and that
$c_{1}\approx 2\times c_{1}^{\text{mirror}}=2/d$ at small $d$---which
exceeds the asymptotic $1/d$ by a factor of $2$ until higher-order
corrections reduce the effective value.  The data instead show
$c_{1}\cdot d\to 1$, suggesting significant cancellation among the
background terms.


% ============================================================================
\begin{thebibliography}{99}

\bibitem{kim2026}
Kim, Y.-H. (2026).
\textit{Phase Quantisation from Optical Lattices to the Zeta Landscape:
A Unified Framework}.
Preprint, arXiv:2026.xxxxx.

\bibitem{titchmarsh1986}
Titchmarsh, E.~C. (1986).
\textit{The Theory of the Riemann Zeta-Function} (2nd ed., revised by
D.~R.~Heath-Brown).
Oxford University Press.

\bibitem{edwards1974}
Edwards, H.~M. (1974).
\textit{Riemann's Zeta Function}.
Academic Press.

\bibitem{davenport1936}
Davenport, H., \& Heilbronn, H. (1936).
On the zeros of certain Dirichlet series.
\textit{Journal of the London Mathematical Society}, 11, 181--185.

\bibitem{silverman1992}
Silverman, J.~H. (1992).
\textit{The Arithmetic of Elliptic Curves} (2nd ed.).
Springer.

\bibitem{iwaniec2004}
Iwaniec, H., \& Kowalski, E. (2004).
\textit{Analytic Number Theory}.
American Mathematical Society.

\bibitem{epstein1903}
Epstein, P. (1903).
Zur Theorie allgemeiner Zetafunktionen.
\textit{Mathematische Annalen}, 56, 615--644.

\bibitem{rubinstein2001}
Rubinstein, M. (2001).
Numerical computations concerning the Riemann hypothesis.
\textit{PhD thesis}, Princeton University.

\bibitem{cremona2019}
Cremona, J.~E. (2019).
\textit{Algorithms for Modular Elliptic Curves} (2nd ed.).
Cambridge University Press.

\bibitem{pari2023}
PARI/GP (2023). Version 2.15.4.
\url{https://pari.math.u-bordeaux.fr/}.

\bibitem{mpmath2023}
Johansson, F., et al.\ (2023).
\textit{mpmath: a Python library for arbitrary-precision
floating-point arithmetic} (version 1.3.0).
\url{https://mpmath.org/}.

\end{thebibliography}

\end{document}
